\chapter{Temporal Trajectories and Segal--Decomposition Execution Semantics}
\label{ch:temporal-trajectories-segal-decomposition-execution}

\section{Temporal orientation and recalled execution semantics}
\label{sec:temporal-orientation-recalled-execution}

\subsection{Automata-theoretic motivation}
\label{subsec:temporal-motivation}

The preceding chapters developed the state semantics of categorical automata.
A state of a Mealy morphism determines a residual behaviour, behavioural
equivalence identifies states with the same residual behaviour, and semantic
factorization produces canonical separated realizations. Those constructions
answer the state-oriented question
\[
    \text{what behaviour is represented by the present state?}
\]

The present chapter studies the complementary execution-oriented question
\[
    \text{what finite run is forced by a state and a finite input path?}
\]

The required one-step structure has already been constructed. A Mealy
morphism
\[
    \Phi=(P,\delta_P,\omega_P):
    \mathbb A\rightsquigarrow\mathbb B
\]
determines an internal execution category whose arrows are oriented in the
direction in which states evolve. Its output operation is equivalently an
internal functor from that execution category to
\(\mathbb B^{\mathrm{op}}\).

The present chapter develops four additional layers of structure.
First, the input labelling functor is shown to be an internal discrete
opfibration, and the execution category is identified as the internal
Grothendieck construction of the deterministic state action. This gives a
representation theorem for Mealy morphisms as two-sided labelled execution
spans.

Second, the construction is extended to all Mealy simulation \(2\)-cells.
A simulation induces a functor between execution categories over the input
interface and a natural transformation comparing their output functors.
Thus the full bicategorical structure of Mealy morphisms is retained.

Third, finite execution paths are organized as sheaves on a site of finite
discrete intervals. These trajectory sheaves form a Grothendieck topos, and
two-sided labelled trajectories form a slice Grothendieck topos. The internal
logic of this topos supplies existential and universal temporal modalities.

Fourth, the augmented nerve supplies coherent identity insertion and
contraction through its strict Segal structure, while its decomposition-space
structure records coherent factorizations of executions. The resulting
objective incidence comonoid remembers all possible cuts of an execution,
rather than only its total composite.

The principal results are therefore:
\begin{enumerate}
    \item every Mealy morphism determines a two-sided labelled execution span
    over \(\mathbb A^{\mathrm{op}}\) and
    \(\mathbb B^{\mathrm{op}}\), whose input leg is an internal discrete
    opfibration;
    \item Mealy morphisms are reconstructed from such labelled discrete-
    opfibration spans;
    \item every Mealy simulation induces an input-strict execution functor
    together with an output-comparison natural transformation;
    \item finite executions are uniquely determined by their initial states
    and finite input paths;
    \item execution categories and finite trajectory objects compose by
    pullback over their common labelled interfaces;
    \item internal graphs are equivalent to trajectory sheaves, and the
    category of trajectory sheaves is a Grothendieck topos;
    \item two-sided labelled trajectories form a slice Grothendieck topos;
    \item internal category structures are equivalent to coherent strict
    Segal composition structures on trajectory sheaves;
    \item execution nerves are decomposition objects and therefore determine
    objective incidence comonoids of execution factorizations;
    \item Mealy simulations induce conservative unique-lifting-of-
    factorization maps and hence morphisms of execution incidence comonoids;
    \item regular semantic quotients induce regular epimorphisms in every
    finite path length and preserve Segal and decomposition structure;
    \item the internal logic of the trajectory topos supplies a typed
    temporal modal calculus;
    \item derivative generation is forward endpoint reachability, while the
    derivative safety core is universal finite-path safety;
    \item in the Boolean specialization, the full word-labelled execution
    category is the free categorical completion of the letterwise derivative
    graph, and its incidence comultiplication is word deconcatenation.
\end{enumerate}

\subsection{Standing assumptions and notation}
\label{subsec:temporal-standing-assumptions}

Throughout the chapter, \(\mathcal E\) is a Grothendieck topos with chosen
pullbacks. All external indexing categories are taken in a fixed universe.
The finite-path, nerve, and decomposition constructions use only finite
limits. The stronger ambient condition is retained because the semantic and
logical applications use regular images, pullback-stable regular
epimorphisms, small coproducts, complete subobject lattices, dependent sums,
dependent products, and Beck--Chevalley isomorphisms.

Write
\[
    \mathbb A=(A_0,A_1,s_A,t_A,e_A,m_A),
    \qquad
    \mathbb B=(B_0,B_1,s_B,t_B,e_B,m_B).
\]
The carrier of \(\Phi\) has anchors
\[
    p_P:P\to A_0,
    \qquad
    q_P:P\to B_0.
\]
Its transition domain is
\[
    D_P
    :=
    A_1\times_{t_A,A_0,p_P}P,
\]
and its structure maps are
\[
    \delta_P:D_P\to P,
    \qquad
    \omega_P:D_P\to B_1.
\]

For
\[
    a:u\to v
\]
and \(x\in P\) with \(p_P(x)=v\), the transition and output comparison
have types
\[
    \delta_P(a,x)\in P_u,
\]
\[
    \omega_P(a,x):
    q_P\bigl(\delta_P(a,x)\bigr)
    \longrightarrow
    q_P(x).
\]
For composable arrows
\[
    a:u\to v,
    \qquad
    b:v\to w,
\]
and \(x\) over \(w\), the Mealy multiplication laws are
\[
    \delta_P(b\circ a,x)
    =
    \delta_P\bigl(a,\delta_P(b,x)\bigr),
\]
\[
    \omega_P(b\circ a,x)
    =
    \omega_P(b,x)
    \circ
    \omega_P\bigl(a,\delta_P(b,x)\bigr).
\]
These equations determine the orientation of every construction below.

\subsection{Recalled execution and output functor}
\label{subsec:recalled-execution-output}

Recall that the execution category of \(\Phi\) has object object and arrow
object
\[
    \operatorname{Exec}(\Phi)_0:=P,
    \qquad
    \operatorname{Exec}(\Phi)_1:=D_P.
\]
Its source and target maps are
\[
    s_{\operatorname{Exec}}(a,x):=x,
    \qquad
    t_{\operatorname{Exec}}(a,x):=\delta_P(a,x).
\]
Thus an execution arrow is
\[
\begin{tikzcd}[column sep=huge]
    x
        \arrow[r, "{(a,x)}"]
    &
    \delta_P(a,x).
\end{tikzcd}
\]

The identity at \(x\) is
\[
    \bigl(e_A(p_Px),x\bigr).
\]
If
\[
    a:u\to v,
    \qquad
    b:v\to w,
\]
and \(x\) lies over \(w\), then
\[
\begin{tikzcd}[column sep=large]
    x
        \arrow[r, "{(b,x)}"]
    &
    \delta_P(b,x)
        \arrow[r, "{(a,\delta_P(b,x))}"]
    &
    \delta_P\bigl(a,\delta_P(b,x)\bigr)
\end{tikzcd}
\]
has composite
\[
\begin{tikzcd}[column sep=huge]
    x
        \arrow[r, "{(b\circ a,x)}"]
    &
    \delta_P(b\circ a,x).
\end{tikzcd}
\]

The output operation is equivalently an internal functor
\[
    O_\Phi:
    \operatorname{Exec}(\Phi)
    \longrightarrow
    \mathbb B^{\mathrm{op}}.
\]
Its object part is \(q_P\), and its arrow part is
\[
    O_{\Phi,1}(a,x)
    =
    \omega_P(a,x)^{\mathrm{op}}.
\]

\subsection{The input labelling functor}
\label{subsec:input-labelling-functor}

The execution category also carries its processed input arrow as a label.

\begin{definition}[Input labelling functor]
\label{def:input-labelling-functor}
The \textbf{input labelling functor}
\[
    I_\Phi:
    \operatorname{Exec}(\Phi)
    \longrightarrow
    \mathbb A^{\mathrm{op}}
\]
has object part
\[
    p_P:P\to A_0
\]
and arrow part
\[
    I_{\Phi,1}(a,x):=a^{\mathrm{op}}.
\]
\end{definition}

For \(a:u\to v\), the execution arrow
\[
    (a,x):x\to\delta_P(a,x)
\]
lies over
\[
    a^{\mathrm{op}}:v\to u.
\]
The reversal is forced by the target-to-source processing convention.

\begin{proposition}[Functoriality of the input label]
\label{prop:functoriality-input-label}
The displayed maps define an internal functor
\[
    I_\Phi:
    \operatorname{Exec}(\Phi)\to\mathbb A^{\mathrm{op}}.
\]
\end{proposition}

\begin{proof}
First, the arrow map preserves source and target. Indeed,
\[
    s_{\mathbb A^{\mathrm{op}}}(a^{\mathrm{op}})
    =
    t_A(a)
    =
    p_P(x)
    =
    p_Ps_{\operatorname{Exec}}(a,x),
\]
and
\[
    t_{\mathbb A^{\mathrm{op}}}(a^{\mathrm{op}})
    =
    s_A(a)
    =
    p_P\bigl(\delta_P(a,x)\bigr)
    =
    p_Pt_{\operatorname{Exec}}(a,x).
\]
The input label of the identity execution at \(x\) is
\[
    e_A(p_Px)^{\mathrm{op}},
\]
which is the identity of \(p_Px\) in
\(\mathbb A^{\mathrm{op}}\).

For a composable pair processed first by \(b\) and then by \(a\), the
execution composite is labelled by
\[
    (b\circ a)^{\mathrm{op}}
    =
    a^{\mathrm{op}}\circ b^{\mathrm{op}},
\]
which is precisely the composite in
\(\mathbb A^{\mathrm{op}}\) in execution order.
\end{proof}

\begin{theorem}[Two-sided labelled execution]
\label{thm:two-sided-labelled-execution}
Every Mealy morphism
\[
    \Phi:\mathbb A\rightsquigarrow\mathbb B
\]
determines a span of internal categories
\[
\begin{tikzcd}[column sep=large]
    \mathbb A^{\mathrm{op}}
    &
    \operatorname{Exec}(\Phi)
        \arrow[l, "I_\Phi"']
        \arrow[r, "O_\Phi"]
    &
    \mathbb B^{\mathrm{op}}.
\end{tikzcd}
\]
Its object map records the input and output anchors of a state, while its
arrow map records the input and output-comparison labels of a one-step
execution.
\end{theorem}

\section{Execution as an internal Grothendieck construction}
\label{sec:execution-grothendieck-construction}

\subsection{Internal discrete opfibrations}
\label{subsec:internal-discrete-opfibrations}

\begin{definition}[Internal discrete opfibration]
\label{def:internal-discrete-opfibration}
An internal functor
\[
    F:\mathbb E\to\mathbb C
\]
is an \textbf{internal discrete opfibration} if the square
\[
\begin{tikzcd}
    E_1
        \arrow[r,"F_1"]
        \arrow[d,"s_E"']
    &
    C_1
        \arrow[d,"s_C"]
    \\
    E_0
        \arrow[r,"F_0"']
    &
    C_0
\end{tikzcd}
\]
is a pullback.
\end{definition}

Internally, this says that for every \(x\in E_0\) and every arrow
\[
    c:F_0(x)\to y
\]
in \(\mathbb C\), there is a unique arrow \(\widetilde c\) in
\(\mathbb E\) with source \(x\) and \(F_1(\widetilde c)=c\).
Its target is then the transported object over \(y\).

Internal discrete opfibrations are stable under pullback and under
composition. Both facts follow directly from the corresponding stability of
pullback squares.

\subsection{The execution lifting theorem}
\label{subsec:execution-lifting-theorem}

\begin{theorem}[Execution lifting]
\label{thm:execution-lifting}
The input labelling functor
\[
    I_\Phi:
    \operatorname{Exec}(\Phi)
    \longrightarrow
    \mathbb A^{\mathrm{op}}
\]
is an internal discrete opfibration.
\end{theorem}

\begin{proof}
The source map of \(\mathbb A^{\mathrm{op}}\) is \(t_A\). Hence the
discrete-opfibration square is
\[
\begin{tikzcd}
    D_P
        \arrow[r,"\pi_{A_1}"]
        \arrow[d,"\pi_P"']
    &
    A_1
        \arrow[d,"t_A"]
    \\
    P
        \arrow[r,"p_P"']
    &
    A_0.
\end{tikzcd}
\]
This is exactly the defining pullback
\[
    D_P=A_1\times_{t_A,A_0,p_P}P.
\]
\end{proof}

\begin{corollary}[Unique one-step lifting]
\label{cor:unique-one-step-lifting}
Let \(x\in P\), and let
\[
    a^{\mathrm{op}}:p_P(x)\to u
\]
be an arrow of \(\mathbb A^{\mathrm{op}}\). There is a unique execution
arrow with source \(x\) and input label \(a^{\mathrm{op}}\), namely
\[
    (a,x):x\to\delta_P(a,x).
\]
\end{corollary}

\subsection{The state action and its Grothendieck construction}
\label{subsec:state-action-grothendieck}

The transition part of \(\Phi\) is equivalently a covariant internal action
of \(\mathbb A^{\mathrm{op}}\) on the object \(p_P:P\to A_0\). On an
arrow
\[
    a^{\mathrm{op}}:v\to u
\]
and a state \(x\in P_v\), the action is
\[
    x\longmapsto\delta_P(a,x)\in P_u.
\]
The Mealy unit and multiplication laws for \(\delta_P\) are exactly the
identity and associativity laws for this action.

\begin{theorem}[Execution as an internal Grothendieck construction]
\label{thm:execution-grothendieck-construction}
The internal category
\[
    \operatorname{Exec}(\Phi)
\]
together with
\[
    I_\Phi:\operatorname{Exec}(\Phi)\to\mathbb A^{\mathrm{op}}
\]
is canonically the internal Grothendieck construction of the state action of
\(\mathbb A^{\mathrm{op}}\) on \(P\).
\end{theorem}

\begin{proof}
The Grothendieck construction of the action has object object \(P\). Its
arrow object consists of pairs
\[
    (a^{\mathrm{op}},x)
\]
with source of \(a^{\mathrm{op}}\) equal to \(p_P(x)\). Since the source
of \(a^{\mathrm{op}}\) is \(t_A(a)\), this arrow object is
\[
    A_1\times_{t_A,A_0,p_P}P=D_P.
\]
Its source is \(x\), its target is the transported state
\(\delta_P(a,x)\), and its composition is determined by the action
multiplication law. These are precisely the object, arrow, source, target,
identity, and composition maps of \(\operatorname{Exec}(\Phi)\).
\end{proof}

\subsection{Reconstruction from a labelled discrete-opfibration span}
\label{subsec:reconstruction-labelled-dopf-span}

Consider a span of internal categories
\[
\begin{tikzcd}[column sep=large]
    \mathbb A^{\mathrm{op}}
    &
    \mathbb E
        \arrow[l,"I"']
        \arrow[r,"O"]
    &
    \mathbb B^{\mathrm{op}},
\end{tikzcd}
\]
where \(I\) is an internal discrete opfibration.

Put
\[
    P:=E_0,
    \qquad
    p:=I_0,
    \qquad
    q:=O_0.
\]
For \((a,x)\in A_1\times_{t_A,A_0,p}P\), let
\[
    \widetilde{a^{\mathrm{op}}}_x
\]
be the unique \(I\)-lift of \(a^{\mathrm{op}}\) with source \(x\).
Define
\[
    \delta(a,x)
    :=
    t_E\bigl(\widetilde{a^{\mathrm{op}}}_x\bigr),
\]
and define \(\omega(a,x)\) by
\[
    O_1\bigl(\widetilde{a^{\mathrm{op}}}_x\bigr)
    =
    \omega(a,x)^{\mathrm{op}}.
\]

\begin{theorem}[Reconstruction of Mealy structure]
\label{thm:reconstruction-mealy-from-dopf-span}
The displayed data define a Mealy morphism
\[
    \Phi_{I,O}:\mathbb A\rightsquigarrow\mathbb B.
\]
Its execution span is canonically isomorphic to
\[
\begin{tikzcd}[column sep=large]
    \mathbb A^{\mathrm{op}}
    &
    \mathbb E
        \arrow[l,"I"']
        \arrow[r,"O"]
    &
    \mathbb B^{\mathrm{op}}.
\end{tikzcd}
\]
Conversely, reconstructing from the execution span of a Mealy morphism
returns its original transition and output maps.
\end{theorem}

\begin{proof}
The target of the unique lift lies over the target of \(a^{\mathrm{op}}\),
which is \(s_A(a)\). Hence
\[
    p\delta(a,x)=s_A(a).
\]
Since \(O\) is a functor, the arrow
\[
    O_1\bigl(\widetilde{a^{\mathrm{op}}}_x\bigr)
\]
has source \(q(x)\) and target \(q(\delta(a,x))\) in
\(\mathbb B^{\mathrm{op}}\). Equivalently,
\[
    \omega(a,x):q(\delta(a,x))\to q(x)
\]
in \(\mathbb B\).

The unique lift of an identity is an identity, so the Mealy unit laws follow
from functoriality of \(\mathbb E\) and \(O\). The composite of the unique
lift of \(b^{\mathrm{op}}\) at \(x\) with the unique lift of
\(a^{\mathrm{op}}\) at its target is the unique lift of
\[
    a^{\mathrm{op}}\circ b^{\mathrm{op}}
    =
    (b\circ a)^{\mathrm{op}}.
\]
The transition multiplication law follows by comparing targets. Applying
\(O\) and passing from \(\mathbb B^{\mathrm{op}}\) to \(\mathbb B\)
gives the output multiplication law.

The arrow object of the reconstructed execution category is the pullback
classifying the unique lifts, hence is canonically \(E_1\). Under this
identification, all internal-category and label maps agree. The converse is
immediate from the explicit lift
\((a,x)\) in the original execution category.
\end{proof}

\begin{corollary}[One-cell execution representation]
\label{cor:one-cell-execution-representation}
For fixed \(\mathbb A\) and \(\mathbb B\), Mealy morphisms
\[
    \mathbb A\rightsquigarrow\mathbb B
\]
are equivalent, up to canonical isomorphism, to two-sided labelled execution
spans
\[
    \mathbb A^{\mathrm{op}}
    \xleftarrow{\ I\ }
    \mathbb E
    \xrightarrow{\ O\ }
    \mathbb B^{\mathrm{op}}
\]
whose input leg \(I\) is an internal discrete opfibration.
\end{corollary}

\subsection{Execution of a pipeline}
\label{subsec:execution-pipeline}

Let
\[
    \Phi=(P,\delta_P,\omega_P):
    \mathbb A\rightsquigarrow\mathbb B,
\]
\[
    \Psi=(Q,\delta_Q,\omega_Q):
    \mathbb B\rightsquigarrow\mathbb C
\]
be composable Mealy morphisms. Their pipeline composite has carrier
\[
    P\times_{q_P,B_0,p_Q}Q.
\]
The common interface of the two execution categories is expressed by
\[
    O_\Phi:
    \operatorname{Exec}(\Phi)\to\mathbb B^{\mathrm{op}},
\]
\[
    I_\Psi:
    \operatorname{Exec}(\Psi)\to\mathbb B^{\mathrm{op}}.
\]

\begin{theorem}[Execution of a pipeline as a pullback]
\label{thm:execution-pipeline-pullback}
There is a canonical isomorphism of internal categories
\[
    \operatorname{Exec}(\Psi\circ\Phi)
    \cong
    \operatorname{Exec}(\Phi)
    \times_{\mathbb B^{\mathrm{op}}}
    \operatorname{Exec}(\Psi).
\]
Under this isomorphism, the input label is induced by \(I_\Phi\), the
output label is induced by \(O_\Psi\), and the resulting input projection
is again an internal discrete opfibration.
\end{theorem}

\begin{proof}
On objects, the pullback is
\[
    P\times_{B_0}Q,
\]
which is the carrier of the pipeline composite.

On arrows, define
\[
    (a,x,y)
    \longmapsto
    \left(
        (a,x),
        (\omega_P(a,x),y)
    \right).
\]
The second execution arrow is defined because
\[
    t_B(\omega_P(a,x))
    =
    q_P(x)
    =
    p_Q(y).
\]
The two arrows have the same label in
\(\mathbb B^{\mathrm{op}}\):
\[
    O_\Phi(a,x)
    =
    \omega_P(a,x)^{\mathrm{op}}
    =
    I_\Psi(\omega_P(a,x),y).
\]

Conversely, an arrow of the pullback consists of
\[
    (a,x)\in\operatorname{Exec}(\Phi)_1,
    \qquad
    (b,y)\in\operatorname{Exec}(\Psi)_1
\]
with
\[
    b^{\mathrm{op}}
    =
    \omega_P(a,x)^{\mathrm{op}}.
\]
Hence \(b=\omega_P(a,x)\), so the data are uniquely determined by
\((a,x,y)\).

For composition, consider a pipeline execution processed first by \(b\) and
then by \(a\). Its two images in the pullback compose to
\[
    \left(
        (b\circ a,x),
        \bigl(
            \omega_P(b,x)
            \circ
            \omega_P(a,\delta_P(b,x)),
            y
        \bigr)
    \right).
\]
The output multiplication law of \(P\) identifies the second input label
with \(\omega_P(b\circ a,x)\). The transition multiplication law of \(Q\)
then identifies the resulting target with the pipeline transition by
\(b\circ a\). Identity preservation is the corresponding unit calculation.
Thus the comparisons are mutually inverse internal functors.

Finally, the projection from the pullback execution category to
\(\operatorname{Exec}(\Phi)\) is a pullback of the discrete opfibration
\(I_\Psi\). Its composite with \(I_\Phi\) is therefore a discrete
opfibration.
\end{proof}

\subsection{Units and coherent associativity}
\label{subsec:execution-units-coherence}

Let
\[
    1_{\mathbb A}:\mathbb A\rightsquigarrow\mathbb A
\]
be the identity Mealy morphism.

\begin{proposition}[Execution unit]
\label{prop:execution-unit}
There is a canonical isomorphism
\[
    \operatorname{Exec}(1_{\mathbb A})
    \cong
    \mathbb A^{\mathrm{op}},
\]
compatible with both input and output labelling.
\end{proposition}

\begin{proof}
The object object is \(A_0\), and
\[
    A_1\times_{t_A,A_0}A_0
    \cong
    A_1.
\]
Under this identification, the execution source and target are \(t_A\) and
\(s_A\), respectively. Both labels send the execution represented by \(a\)
to \(a^{\mathrm{op}}\).
\end{proof}

\begin{proposition}[Coherence of execution pullbacks]
\label{prop:coherence-execution-pullbacks}
The pipeline comparison isomorphisms are compatible with the associativity
and unit isomorphisms of Mealy composition.
\end{proposition}

\begin{proof}
For three composable Mealy morphisms, both parenthesized pullbacks of
execution categories are presentations of the same iterated pullback over
the two common interfaces. Every route around the associativity comparison
has the same projections to the three execution-category factors. Equality
therefore follows from the pullback universal property. The unit comparison
is analogous.
\end{proof}

\section{Simulation \(2\)-cells and bicategorical execution semantics}
\label{sec:simulation-two-cells-execution}

The preceding representation theorem concerns Mealy morphisms as
\(1\)-cells. We now retain the full simulation \(2\)-cells of the established
Mealy bicategory.

\subsection{From simulations to execution functors}
\label{subsec:simulations-execution-functors}

Let
\[
    (\theta,\alpha):
    \Phi\Longrightarrow\Psi
\]
be a Mealy simulation between
\[
    \Phi=(P,\delta_P,\omega_P),
    \qquad
    \Psi=(Q,\delta_Q,\omega_Q):
    \mathbb A\rightsquigarrow\mathbb B.
\]
With the established displayed orientation, its state map has direction
\[
    \theta:Q\to P.
\]
It lies over the input anchor and is transition equivariant:
\[
    p_P\theta=p_Q,
\]
\[
    \theta\bigl(\delta_Q(a,y)\bigr)
    =
    \delta_P(a,\theta y).
\]
Its output-comparison component has type
\[
    \alpha_y:
    q_P(\theta y)
    \longrightarrow
    q_Q(y)
\]
in \(\mathbb B\), and satisfies
\[
    \alpha_y
    \circ
    \omega_P(a,\theta y)
    =
    \omega_Q(a,y)
    \circ
    \alpha_{\delta_Q(a,y)}.
\]

\begin{proposition}[Execution functor of a simulation]
\label{prop:execution-functor-simulation}
The state map of a Mealy simulation induces an internal functor
\[
    \operatorname{Exec}(\theta):
    \operatorname{Exec}(\Psi)
    \longrightarrow
    \operatorname{Exec}(\Phi)
\]
defined by
\[
    y\longmapsto\theta y,
\]
\[
    (a,y)\longmapsto(a,\theta y).
\]
It lies strictly over the input interface:
\[
    I_\Phi\operatorname{Exec}(\theta)=I_\Psi.
\]
\end{proposition}

\begin{proof}
The source is plainly preserved. Transition equivariance gives
\[
    \theta\bigl(\delta_Q(a,y)\bigr)
    =
    \delta_P(a,\theta y),
\]
so targets are preserved. The input arrow is unchanged, and therefore
identities and composition are preserved by the corresponding formulas in
the two execution categories. The equality over
\(\mathbb A^{\mathrm{op}}\) follows on objects from
\(p_P\theta=p_Q\) and on arrows from preservation of the input label.
\end{proof}

\subsection{Output comparison as a natural transformation}
\label{subsec:simulation-output-natural-transformation}

Passing the simulation component to the opposite output category gives
\[
    \alpha_y^{\mathrm{op}}:
    q_Q(y)
    \longrightarrow
    q_P(\theta y)
\]
in \(\mathbb B^{\mathrm{op}}\).

\begin{proposition}[Output naturality of a simulation]
\label{prop:simulation-output-naturality}
The components \(\alpha_y^{\mathrm{op}}\) define an internal natural
transformation
\[
    \alpha^{\mathrm{op}}:
    O_\Psi
    \Longrightarrow
    O_\Phi\operatorname{Exec}(\theta).
\]
\end{proposition}

\begin{proof}
For an execution arrow
\[
    (a,y):y\to\delta_Q(a,y),
\]
naturality in \(\mathbb B^{\mathrm{op}}\) is the equality
\[
    \omega_P(a,\theta y)^{\mathrm{op}}
    \circ
    \alpha_y^{\mathrm{op}}
    =
    \alpha_{\delta_Q(a,y)}^{\mathrm{op}}
    \circ
    \omega_Q(a,y)^{\mathrm{op}}.
\]
Passing to \(\mathbb B\) gives
\[
    \alpha_y
    \circ
    \omega_P(a,\theta y)
    =
    \omega_Q(a,y)
    \circ
    \alpha_{\delta_Q(a,y)},
\]
which is exactly the simulation output law.
\end{proof}

\subsection{The bicategory of output-lax discrete-opfibration spans}
\label{subsec:bicategory-output-lax-dopf-spans}

\begin{definition}[Output-lax discrete-opfibration span]
\label{def:output-lax-dopf-span}
An \textbf{output-lax discrete-opfibration span} from an internal category
\(\mathbb C\) to an internal category \(\mathbb D\) is a span
\[
\begin{tikzcd}[column sep=large]
    \mathbb C
    &
    \mathbb E
        \arrow[l,"I"']
        \arrow[r,"O"]
    &
    \mathbb D
\end{tikzcd}
\]
in which \(I\) is an internal discrete opfibration.
\end{definition}

A \(2\)-cell displayed from
\[
    (I_P,O_P):\mathbb C\leftarrow\mathbb E_P\to\mathbb D
\]
to
\[
    (I_Q,O_Q):\mathbb C\leftarrow\mathbb E_Q\to\mathbb D
\]
consists of:
\begin{enumerate}
    \item an internal functor
    \[
        F:\mathbb E_Q\to\mathbb E_P
    \]
    satisfying
    \[
        I_PF=I_Q;
    \]
    \item an internal natural transformation
    \[
        \gamma:
        O_Q\Longrightarrow O_PF.
    \]
\end{enumerate}
The reversal in the apex-functor direction matches the displayed direction of
Mealy simulations.

Vertical composition is defined by composition of the apex functors and
pasting of the output transformations. Horizontal composition is induced by
pullback over the common interface, together with whiskering and pasting of
the output transformations.

\begin{proposition}[Bicategory of labelled execution spans]
\label{prop:bicategory-labelled-execution-spans}
Internal categories, output-lax discrete-opfibration spans, and the displayed
\(2\)-cells form a bicategory
\[
    \mathsf{DOpSpan}^{\mathrm{out}}(\mathcal E).
\]
Its associators and unitors are induced by the canonical pullback
isomorphisms.
\end{proposition}

\begin{proof}
Composition of \(1\)-cells is pullback composition of spans. If
\[
    \mathbb C\xleftarrow{I}\mathbb E\xrightarrow{O}\mathbb D
\]
and
\[
    \mathbb D\xleftarrow{J}\mathbb F\xrightarrow{R}\mathbb K
\]
are composable, the composite apex is
\[
    \mathbb E\times_{\mathbb D}\mathbb F.
\]
The projection to \(\mathbb E\) is a pullback of the discrete opfibration
\(J\), hence is a discrete opfibration. Its composite with \(I\) is
therefore a discrete opfibration.

Vertical and horizontal composition of \(2\)-cells follow from ordinary
composition, pullback functoriality, and pasting of natural transformations.
The interchange law follows from functoriality of these constructions.
Associativity and unit coherence reduce to equality of pullback-induced maps
with the same projections.
\end{proof}

\subsection{The bicategorical execution representation}
\label{subsec:bicategorical-execution-representation}

\begin{theorem}[Bicategorical Grothendieck execution representation]
\label{thm:bicategorical-grothendieck-execution}
The assignments
\[
    \mathbb A\longmapsto\mathbb A^{\mathrm{op}},
\]
\[
    \Phi\longmapsto
    \left(
        \mathbb A^{\mathrm{op}}
        \xleftarrow{I_\Phi}
        \operatorname{Exec}(\Phi)
        \xrightarrow{O_\Phi}
        \mathbb B^{\mathrm{op}}
    \right),
\]
and
\[
    (\theta,\alpha)
    \longmapsto
    \left(
        \operatorname{Exec}(\theta),
        \alpha^{\mathrm{op}}
    \right)
\]
define a pseudofunctor
\[
    \operatorname{Mly}(\mathcal E)
    \longrightarrow
    \mathsf{DOpSpan}^{\mathrm{out}}(\mathcal E).
\]
It is locally an equivalence onto the full sub-bicategory of labelled
execution spans whose left legs are internal discrete opfibrations.
Consequently, the execution construction gives a biequivalent presentation
of Mealy morphisms and Mealy simulations by labelled Grothendieck execution
spans.
\end{theorem}

\begin{proof}
The \(1\)-cell comparison is
Theorem~\ref{thm:reconstruction-mealy-from-dopf-span}. Propositions
\ref{prop:execution-functor-simulation} and
\ref{prop:simulation-output-naturality} give the action on \(2\)-cells.
Vertical composition agrees because state maps compose and output comparison
arrows paste in the same order as natural transformations.

For horizontal composition, the execution category of a pipeline is the
pullback of the two execution categories. The horizontal composite of state
maps is the pullback-induced map between these pullbacks, and the horizontal
simulation output comparison is exactly the pasted comparison obtained after
whiskering along the two pipeline factors. Hence the pipeline comparison
isomorphisms define the pseudofunctorial compositor. The execution unit gives
the unitor.

Conversely, let
\[
    F:\mathbb E_Q\to\mathbb E_P
\]
be a functor over the input interface, together with
\[
    \gamma:O_Q\Rightarrow O_PF.
\]
The object map of \(F\) gives a state map \(\theta:Q\to P\). Since both
input legs are discrete opfibrations and \(F\) preserves input labels,
its arrow map is forced to send the unique lift at \(y\) to the unique lift
at \(\theta y\). This is precisely transition equivariance. Passing
\(\gamma\) to \(\mathbb B\) gives the simulation output component, and
naturality gives the simulation equation. These constructions are mutually
inverse up to the canonical execution-span isomorphisms.
\end{proof}

\subsection{Strict semantic homomorphisms}
\label{subsec:execution-strict-homomorphisms}

Let
\[
    f:P\to Q
\]
be a strict semantic homomorphism between Mealy morphisms over the same
interfaces. It may be regarded as a strict Mealy simulation displayed in the
reverse direction
\[
    Q\Longrightarrow P.
\]

\begin{proposition}[Strict functoriality of execution]
\label{prop:functoriality-execution}
The map \(f\) induces an internal functor
\[
    \operatorname{Exec}(f):
    \operatorname{Exec}(P)
    \longrightarrow
    \operatorname{Exec}(Q)
\]
by
\[
    x\longmapsto f(x),
    \qquad
    (a,x)\longmapsto(a,f(x)),
\]
satisfying
\[
    I_Q\operatorname{Exec}(f)=I_P,
    \qquad
    O_Q\operatorname{Exec}(f)=O_P.
\]
\end{proposition}

\begin{proof}
Strict transition preservation gives
\[
    f\bigl(\delta_P(a,x)\bigr)
    =
    \delta_Q(a,f(x)),
\]
and strict output preservation gives
\[
    \omega_Q(a,f(x))
    =
    \omega_P(a,x).
\]
The remaining assertions follow directly.
\end{proof}

The distinction between general simulations and strict semantic
homomorphisms will remain visible below. General simulations act
contravariantly on execution apexes and compare output labels by natural
transformation. Strict homomorphisms induce actual forward morphisms of
strictly labelled trajectory objects and are therefore the maps used for
semantic quotients.

\section{Finite labelled executions}
\label{sec:finite-labelled-executions}

\subsection{Finite execution objects}
\label{subsec:finite-execution-objects}

For \(n\geq1\), define
\[
    \mathscr T_\Phi(n)
    :=
    \underbrace{
    \operatorname{Exec}(\Phi)_1
    \times_P
    \cdots
    \times_P
    \operatorname{Exec}(\Phi)_1
    }_{n\text{ factors}},
\]
where the pullbacks identify the target of one execution step with the source
of the next. Put
\[
    \mathscr T_\Phi(0):=P,
    \qquad
    \mathscr T_\Phi(-1):=1.
\]

A generalized element of \(\mathscr T_\Phi(n)\) is written
\[
\begin{tikzcd}[column sep=large]
    x_0
        \arrow[r, "{(a_1,x_0)}"]
    &
    x_1
        \arrow[r, "{(a_2,x_1)}"]
    &
    \cdots
        \arrow[r]
    &
    x_n,
\end{tikzcd}
\]
where
\[
    x_i=\delta_P(a_i,x_{i-1})
\]
and
\[
    a_i:p_P(x_i)\to p_P(x_{i-1})
\]
in \(\mathbb A\).

The corresponding input path in
\(\mathbb A^{\mathrm{op}}\) is
\[
\begin{tikzcd}[column sep=large]
    p_P(x_0)
        \arrow[r, "a_1^{\mathrm{op}}"]
    &
    p_P(x_1)
        \arrow[r, "a_2^{\mathrm{op}}"]
    &
    \cdots
        \arrow[r]
    &
    p_P(x_n),
\end{tikzcd}
\]
and the output-comparison path in
\(\mathbb B^{\mathrm{op}}\) is
\[
\begin{tikzcd}[column sep=large]
    q_P(x_0)
        \arrow[r, "{\omega_P(a_1,x_0)^{\mathrm{op}}}"]
    &
    q_P(x_1)
        \arrow[r, "{\omega_P(a_2,x_1)^{\mathrm{op}}}"]
    &
    \cdots
        \arrow[r]
    &
    q_P(x_n).
\end{tikzcd}
\]

For \(0\leq i\leq n\), write
\[
    \nu_i^{\Phi,n}:
    \mathscr T_\Phi(n)\to P
\]
for the \(i\)-th state map. In particular,
\[
    s_n^\Phi:=\nu_0^{\Phi,n},
    \qquad
    t_n^\Phi:=\nu_n^{\Phi,n}.
\]

For \(0\leq i\leq j\leq n\), projection to the factors between boundaries
\(i\) and \(j\) gives a restriction map
\[
    \operatorname{res}_{i,j}:
    \mathscr T_\Phi(n)
    \longrightarrow
    \mathscr T_\Phi(j-i).
\]
When \(i=j\), this is the vertex map
\(\nu_i^{\Phi,n}\). Restriction to the empty interval is the unique map to
\(1\).

Adjacent finite executions glue along their common boundary state. For
\(m,n\geq0\), the iterated-pullback presentations give a canonical
isomorphism
\[
    \mathscr T_\Phi(m+n)
    \cong
    \mathscr T_\Phi(m)
    \times_P
    \mathscr T_\Phi(n),
\]
where the pullback uses the terminal boundary of the first execution and the
initial boundary of the second.

\subsection{Finite input-path objects}
\label{subsec:finite-input-path-objects}

For \(n\geq2\), define the object of \(n\)-step paths in
\(\mathbb A^{\mathrm{op}}\) by
\[
    \mathbb A^{\mathrm{op}}[n]
    :=
    \underbrace{
    A_1
    \times_{s_A,A_0,t_A}
    A_1
    \times_{s_A,A_0,t_A}
    \cdots
    \times_{s_A,A_0,t_A}
    A_1
    }_{n\text{ factors}}.
\]
Put
\[
    \mathbb A^{\mathrm{op}}[1]:=A_1,
    \qquad
    \mathbb A^{\mathrm{op}}[0]:=A_0,
    \qquad
    \mathbb A^{\mathrm{op}}[-1]:=1.
\]

A generalized element
\[
    (a_1,\ldots,a_n)
\]
satisfies
\[
    s_A(a_i)=t_A(a_{i+1})
\]
and represents the path
\[
\begin{tikzcd}[column sep=large]
    t_A(a_1)
        \arrow[r, "a_1^{\mathrm{op}}"]
    &
    s_A(a_1)
        \arrow[r, "a_2^{\mathrm{op}}"]
    &
    \cdots
        \arrow[r, "a_n^{\mathrm{op}}"]
    &
    s_A(a_n).
\end{tikzcd}
\]

Write
\[
    \lambda_0^{A,n}:
    \mathbb A^{\mathrm{op}}[n]\to A_0
\]
for the initial-boundary map. The input labelling functor induces
\[
    I_{\Phi,n}:
    \mathscr T_\Phi(n)
    \longrightarrow
    \mathbb A^{\mathrm{op}}[n].
\]

\subsection{Deterministic lifting of finite input paths}
\label{subsec:deterministic-path-lifting}

\begin{proposition}[Deterministic path lifting]
\label{prop:deterministic-path-lifting}
For every \(n\geq0\), the canonical map
\[
    \left(
        s_n^\Phi,
        I_{\Phi,n}
    \right):
    \mathscr T_\Phi(n)
    \longrightarrow
    P\times_{A_0}\mathbb A^{\mathrm{op}}[n]
\]
is an isomorphism. The pullback uses
\[
    p_P:P\to A_0
\]
and
\[
    \lambda_0^{A,n}:
    \mathbb A^{\mathrm{op}}[n]\to A_0.
\]
\end{proposition}

\begin{proof}
The nerve of an internal discrete opfibration is cartesian over its object
part in every degree. Applying this to
\[
    I_\Phi:\operatorname{Exec}(\Phi)\to\mathbb A^{\mathrm{op}}
\]
gives the displayed pullback isomorphism.

Explicitly, a finite execution determines its initial state and input path.
Conversely, given an initial state \(x_0\) and a path represented by arrows
\[
    a_i:u_i\to u_{i-1},
    \qquad
    u_0=p_P(x_0),
\]
define recursively
\[
    x_i:=\delta_P(a_i,x_{i-1}).
\]
The transition typing equations give \(p_P(x_i)=u_i\), so every successive
transition is defined. Uniqueness follows from unique lifting for the
discrete opfibration \(I_\Phi\).
\end{proof}

\begin{corollary}[Finite execution determinism]
\label{cor:finite-execution-determinism}
A finite execution is uniquely determined by its initial state and its input
path. Its complete state path and output-comparison path are then forced by
the transition and output operations.
\end{corollary}

\subsection{Functoriality on finite paths}
\label{subsec:finite-path-functoriality}

A strict semantic homomorphism
\[
    f:P\to Q
\]
induces maps
\[
    \mathscr T_f(n):
    \mathscr T_P(n)
    \longrightarrow
    \mathscr T_Q(n)
\]
that apply \(f\) at every state boundary and leave the input and output paths
fixed.

A general simulation
\[
    (\theta,\alpha):P\Longrightarrow Q
\]
induces maps in the opposite apex direction
\[
    \mathscr T_\theta(n):
    \mathscr T_Q(n)
    \longrightarrow
    \mathscr T_P(n),
\]
obtained by applying \(\theta\) at every state boundary. The input path is
preserved strictly. The output comparison \(\alpha\) determines a coherent
boundarywise comparison between the output paths, which will be encoded
simplicially in
Section~\ref{subsec:simulations-simplicial-prisms}.

\section{The finite-path site and the trajectory topos}
\label{sec:finite-path-site-trajectory-topos}

The finite execution objects use only the underlying graph of the execution
category. We first isolate this graph-level temporal structure, then identify
its category of models as a Grothendieck topos.

\subsection{The augmented discrete interval category}
\label{subsec:augmented-discrete-interval-category}

Let
\[
    [-1]:=\varnothing.
\]
For \(n\geq0\), write
\[
    [n]:=\{0<1<\cdots<n\}.
\]
The object \([n]\) represents an interval with \(n\) segments and
\(n+1\) boundary instants.

\begin{definition}[Augmented discrete interval category]
\label{def:augmented-discrete-interval-category}
The \textbf{augmented discrete interval category}
\[
    \mathbf{Int}_{\mathrm{disc}}
\]
has objects
\[
    [-1],[0],[1],[2],\ldots.
\]
A morphism
\[
    \iota:[k]\longrightarrow[n]
\]
is an injective order-preserving map whose image is a contiguous
subinterval. For \(k\geq0\), it is uniquely determined by a starting
position \(i\) and has the form
\[
    \iota(r)=i+r.
\]
There is a unique morphism
\[
    [-1]\to[n]
\]
for every \(n\geq-1\).
\end{definition}

Write
\[
    (i,k):[k]\to[n]
\]
for the inclusion beginning at \(i\). Composition is
\[
    (i,k)(j,\ell)=(i+j,\ell).
\]

The maps
\[
    \epsilon_0,\epsilon_1:[0]\rightrightarrows[1]
\]
select the initial and terminal boundaries. For \(0\leq i<n\), write
\[
    e_i:[1]\to[n]
\]
for the elementary edge inclusion with image \(\{i,i+1\}\).

Intersections of contiguous subintervals are contiguous, possibly empty.
Hence pullbacks of interval inclusions exist in
\(\mathbf{Int}_{\mathrm{disc}}\).

\subsection{The elementary-edge topology}
\label{subsec:elementary-edge-topology}

\begin{definition}[Elementary path covers]
\label{def:elementary-path-covers}
For \(n\geq1\), the \textbf{elementary path cover} of \([n]\) is
\[
    \{e_i:[1]\to[n]\}_{0\leq i<n}.
\]
The identity family covers \([0]\), and the empty family covers \([-1]\).
\end{definition}

\begin{definition}[Finite-path topology]
\label{def:finite-path-topology}
The \textbf{finite-path topology}
\[
    J_{\mathrm{path}}
\]
is the Grothendieck topology generated by the elementary path covers, the
identity cover of \([0]\), and the empty cover of \([-1]\).
\end{definition}

\begin{proposition}[Basis property]
\label{prop:path-covers-form-basis}
The displayed covering families form a basis for
\(J_{\mathrm{path}}\).
\end{proposition}

\begin{proof}
Let \(S_n\) be the sieve generated by the elementary edges of \([n]\), and
let
\[
    \iota:[k]\to[n]
\]
be a contiguous inclusion.

If \(k\geq1\), the pullback sieve \(\iota^\ast S_n\) contains every
elementary edge inclusion of \([k]\). If \(k=0\), it contains the identity
of \([0]\). If \(k=-1\), it is maximal, since the unique map from
\([-1]\) factors through every elementary edge. Thus pullbacks of the
generating covers are covering.

Every member \(e_i:[1]\to[n]\) of an elementary cover is itself covered by
the identity cover of \([1]\). Composing these covers reproduces the
original elementary cover. The specified covers of \([0]\) and \([-1]\)
are respectively the identity and empty families. Hence the basis
transitivity condition holds.
\end{proof}

\subsection{Trajectory sheaves}
\label{subsec:trajectory-sheaves}

\begin{definition}[Trajectory presheaf]
\label{def:trajectory-presheaf}
An \(\mathcal E\)-valued \textbf{trajectory presheaf} is a functor
\[
    X:
    \mathbf{Int}_{\mathrm{disc}}^{\mathrm{op}}
    \to
    \mathcal E.
\]
For a contiguous inclusion
\[
    \iota:[k]\to[n],
\]
the morphism
\[
    X(\iota):X(n)\to X(k)
\]
is called contiguous path restriction.
\end{definition}

\begin{definition}[Trajectory sheaf]
\label{def:trajectory-sheaf}
A trajectory presheaf is a \textbf{trajectory sheaf} if it is a sheaf for
\(J_{\mathrm{path}}\).
\end{definition}

Write
\[
    \mathbf{TrajSh}_{\mathcal E}
    :=
    \mathbf{Sh}_{\mathcal E}
    \left(
        \mathbf{Int}_{\mathrm{disc}},
        J_{\mathrm{path}}
    \right).
\]

\begin{proposition}[Concrete trajectory-sheaf condition]
\label{prop:concrete-trajectory-sheaf-condition}
A trajectory presheaf \(X\) is a trajectory sheaf if and only if:
\begin{enumerate}
    \item the canonical map
    \[
        X(-1)\to1
    \]
    is an isomorphism;
    \item for every \(n\geq1\), the elementary-edge restriction map
    \[
        X(n)
        \longrightarrow
        \underbrace{
        X(1)\times_{X(0)}
        \cdots\times_{X(0)}X(1)
        }_{n\text{ factors}}
    \]
    is an isomorphism.
\end{enumerate}
The pullbacks use terminal restriction on one edge and initial restriction
on the next.
\end{proposition}

\begin{proof}
The matching object for the empty cover of \([-1]\) is terminal, so the
sheaf condition there is exactly
\[
    X(-1)\cong1.
\]

For the elementary cover of \([n]\), a matching family consists of
sections
\[
    x_i\in X(1),
    \qquad
    0\leq i<n,
\]
whose adjacent endpoint restrictions agree:
\[
    X(\epsilon_1)(x_i)
    =
    X(\epsilon_0)(x_{i+1}).
\]
These are precisely the pullback equations defining the displayed iterated
pullback. Compatibility on nonadjacent intersections is automatic because
those intersections are empty and \(X(-1)\cong1\).
\end{proof}

\begin{corollary}[Adjacent temporal gluing]
\label{cor:adjacent-temporal-gluing}
For every trajectory sheaf \(X\) and \(m,n\in\mathbb N\), there is a
canonical isomorphism
\[
    X(m+n)
    \cong
    X(m)\times_{X(0)}X(n),
\]
where the pullback identifies the terminal boundary of the first path with
the initial boundary of the second.
\end{corollary}

\begin{proof}
Both sides are canonically the iterated pullback of \(m+n\) copies of
\(X(1)\) over \(X(0)\).
\end{proof}

\subsection{Path objects of an internal graph}
\label{subsec:path-objects-internal-graph}

Let
\[
    G=
    \left(
    G_1
    \mathrel{\substack{\xrightarrow{t_G}\\[-0.4em]\xrightarrow[s_G]{}}}
    G_0
    \right)
\]
be an internal graph.

Define
\[
    G^{[-1]}:=1,
    \qquad
    G^{[0]}:=G_0,
    \qquad
    G^{[1]}:=G_1,
\]
and, for \(n\geq2\),
\[
    G^{[n]}
    :=
    \underbrace{
    G_1\times_{G_0}\cdots\times_{G_0}G_1
    }_{n\text{ factors}},
\]
where the target of one edge is identified with the source of the next.
A contiguous inclusion selects the corresponding edge block and boundary
vertices.

\begin{definition}[Graph trajectory presheaf]
\label{def:graph-trajectory-presheaf}
The \textbf{graph trajectory presheaf} is
\[
    \operatorname{Path}(G):
    \mathbf{Int}_{\mathrm{disc}}^{\mathrm{op}}
    \to
    \mathcal E,
    \qquad
    \operatorname{Path}(G)(n):=G^{[n]}.
\]
\end{definition}

\begin{proposition}[Temporal gluing of graph paths]
\label{prop:temporal-gluing-graph-paths}
For every internal graph \(G\),
\[
    \operatorname{Path}(G)
\]
is a trajectory sheaf.
\end{proposition}

\begin{proof}
Its value at \([-1]\) is terminal. In every positive degree, the
 elementary-edge comparison is the defining iterated-pullback
identification.
\end{proof}

\subsection{Reconstruction of graphs}
\label{subsec:reconstruction-graphs}

Let \(X\) be a trajectory sheaf. Define an internal graph \(G_X\) by
\[
    (G_X)_0:=X(0),
    \qquad
    (G_X)_1:=X(1),
\]
with
\[
    s_X:=X(\epsilon_0),
    \qquad
    t_X:=X(\epsilon_1).
\]

\begin{proposition}[Reconstruction from degrees zero and one]
\label{prop:reconstruction-degrees-zero-one}
For every trajectory sheaf \(X\), there are canonical isomorphisms
\[
    \theta_{X,n}:
    X(n)
    \xrightarrow{\cong}
    (G_X)^{[n]}
\]
for all \(n\geq-1\), natural in contiguous interval inclusions.
\end{proposition}

\begin{proof}
The cases \(-1,0,1\) follow from the definitions. For \(n\geq2\), apply
the concrete sheaf condition. Under the resulting elementary-edge
isomorphisms, a contiguous inclusion
\[
    (i,k):[k]\to[n]
\]
acts on both sides by projection onto the block of edges
\[
    i,i+1,\ldots,i+k-1
\]
and the corresponding boundary vertices. Hence the isomorphisms are natural.
\end{proof}

\begin{theorem}[Graph--trajectory-sheaf equivalence]
\label{thm:graph-trajectory-sheaf-equivalence}
The graph path construction defines an equivalence
\[
\begin{tikzcd}[column sep=large]
    \mathbf{Graph}_{\mathrm{int}}(\mathcal E)
        \arrow[r, shift left=0.7ex, "{\operatorname{Path}}"]
    &
    \mathbf{TrajSh}_{\mathcal E}
        \arrow[l, shift left=0.7ex, "{X\mapsto G_X}"] .
\end{tikzcd}
\]
\end{theorem}

\begin{proof}
Reconstruction of \(\operatorname{Path}(G)\) in degrees zero and one
returns \(G\) strictly. For a trajectory sheaf \(X\),
Proposition~\ref{prop:reconstruction-degrees-zero-one} gives
\[
    X\cong\operatorname{Path}(G_X).
\]
A graph morphism induces the corresponding maps between all finite iterated
pullbacks. Conversely, a morphism of trajectory sheaves is determined by its
degree-zero and degree-one components, since every higher component is
forced by the elementary-edge isomorphisms.
\end{proof}

\begin{corollary}[No independent higher graph data]
\label{cor:no-independent-higher-graph-data}
A trajectory sheaf contains no independent graph-level data in lengths
greater than one. Every finite path object is reconstructed from
\[
    X(0),
    \qquad
    X(1),
    \qquad
    X(1)\rightrightarrows X(0).
\]
\end{corollary}

Applied to the underlying graph of \(\operatorname{Exec}(\Phi)\), this
gives
\[
    \operatorname{Path}
    \bigl(
        U\operatorname{Exec}(\Phi)
    \bigr)(n)
    \cong
    \mathscr T_\Phi(n).
\]

\subsection{The temporal trajectory topos}
\label{subsec:temporal-trajectory-topos}

Let \(\mathbb G\) denote the walking directed-graph category, with objects
\(0,1\) and two parallel arrows
\[
    s,t:1\rightrightarrows0.
\]
Then
\[
    \mathbf{Graph}_{\mathrm{int}}(\mathcal E)
    \cong
    \mathcal E^{\mathbb G}.
\]

\begin{theorem}[Temporal trajectory topos]
\label{thm:temporal-trajectory-topos}
The category
\[
    \mathbf{TrajSh}_{\mathcal E}
\]
is a Grothendieck topos. More precisely, there are canonical equivalences
\[
    \mathbf{TrajSh}_{\mathcal E}
    \simeq
    \mathbf{Graph}_{\mathrm{int}}(\mathcal E)
    \cong
    \mathcal E^{\mathbb G}.
\]
\end{theorem}

\begin{proof}
The first equivalence is
Theorem~\ref{thm:graph-trajectory-sheaf-equivalence}. Since \(\mathbb G\)
is small and \(\mathcal E\) is a Grothendieck topos, the functor category
\(\mathcal E^{\mathbb G}\) is again a Grothendieck topos.
\end{proof}

Consequently, limits and colimits of trajectory sheaves are computed on
vertex and edge objects. The category is locally cartesian closed, has power
objects and a subobject classifier, and supports dependent sums and dependent
products along every trajectory morphism.

There is an essential geometric morphism
\[
    \pi:
    \mathbf{TrajSh}_{\mathcal E}
    \longrightarrow
    \mathcal E
\]
whose inverse-image functor sends an object \(X\in\mathcal E\) to the
constant trajectory graph
\[
    X
    \mathrel{\substack{\xrightarrow{1_X}\\[-0.4em]\xrightarrow[1_X]{}}}
    X.
\]
Under the graph presentation, its adjoints are
\[
    \pi_!(G)
    =
    \operatorname{coeq}
    \left(
        G_1
        \mathrel{\substack{\xrightarrow{s_G}\\[-0.4em]\xrightarrow[t_G]{}}}
        G_0
    \right),
\]
and
\[
    \pi_*(G)
    =
    \operatorname{eq}
    \left(
        G_1
        \mathrel{\substack{\xrightarrow{s_G}\\[-0.4em]\xrightarrow[t_G]{}}}
        G_0
    \right).
\]
Thus \(\pi_!\) forms graph components, while \(\pi_*\) selects loop edges.

\subsection{The topos of labelled trajectories}
\label{subsec:topos-labelled-trajectories}

Fix internal categories \(\mathbb A\) and \(\mathbb B\). Put
\[
    \mathscr L_{\mathbb A,\mathbb B}
    :=
    \operatorname{Path}(\mathbb A^{\mathrm{op}})
    \times
    \operatorname{Path}(\mathbb B^{\mathrm{op}}).
\]

\begin{definition}[Labelled trajectory topos]
\label{def:labelled-trajectory-topos}
The \textbf{topos of two-sided labelled trajectories} is the slice
\[
    \mathbf{LabTraj}_{\mathbb A,\mathbb B}
    :=
    \mathbf{TrajSh}_{\mathcal E}
    /
    \mathscr L_{\mathbb A,\mathbb B}.
\]
\end{definition}

An object is a trajectory sheaf \(X\) equipped with maps
\[
    L_A:X\to\operatorname{Path}(\mathbb A^{\mathrm{op}}),
\]
\[
    L_B:X\to\operatorname{Path}(\mathbb B^{\mathrm{op}}).
\]
Equivalently, it is a two-sided labelled trajectory span.

\begin{corollary}[Labelled trajectory topos]
\label{cor:labelled-trajectory-topos}
The category
\[
    \mathbf{LabTraj}_{\mathbb A,\mathbb B}
\]
is a Grothendieck topos.
\end{corollary}

\begin{proof}
It is a slice of the Grothendieck topos
\(\mathbf{TrajSh}_{\mathcal E}\).
\end{proof}

The labelled trajectory topos is an ambient universe for deterministic and
nondeterministic trajectory systems, trajectory relations, partial systems,
and input/output predicates. Mealy execution trajectories form a
distinguished class of objects characterized by deterministic input lifting
and by the existence of compatible Segal composition.

\begin{remark}[Strict and lax label comparison]
\label{rem:strict-lax-label-comparison}
Morphisms in the slice topos preserve both labels strictly. Strict semantic
homomorphisms therefore induce morphisms in
\(\mathbf{LabTraj}_{\mathbb A,\mathbb B}\).
General Mealy simulations preserve input labels strictly but compare output
labels by a natural transformation. Their full semantics belongs to the
output-lax bicategorical and simplicial structures, not merely to the
underlying slice category.
\end{remark}

\section{Labelled finite-trajectory semantics}
\label{sec:labelled-trajectory-semantics}

\subsection{The execution trajectory sheaf}
\label{subsec:execution-trajectory-sheaf}

\begin{definition}[Execution trajectory sheaf]
\label{def:execution-trajectory-sheaf}
The \textbf{execution trajectory sheaf} of
\[
    \Phi:\mathbb A\rightsquigarrow\mathbb B
\]
is
\[
    \mathscr T_\Phi
    :=
    \operatorname{Path}
    \bigl(
        \operatorname{Exec}(\Phi)
    \bigr).
\]
\end{definition}

Its degree-\(n\) object is canonically the finite execution object introduced
in Section~\ref{sec:finite-labelled-executions}.

The two label functors induce morphisms of trajectory sheaves
\[
    \operatorname{Path}(I_\Phi):
    \mathscr T_\Phi
    \to
    \operatorname{Path}(\mathbb A^{\mathrm{op}}),
\]
\[
    \operatorname{Path}(O_\Phi):
    \mathscr T_\Phi
    \to
    \operatorname{Path}(\mathbb B^{\mathrm{op}}).
\]

\begin{definition}[Labelled finite-trajectory semantics]
\label{def:labelled-finite-trajectory-semantics}
The \textbf{labelled finite-trajectory semantics} of \(\Phi\) is the object
of \(\mathbf{LabTraj}_{\mathbb A,\mathbb B}\) represented by
\[
\begin{tikzcd}[column sep=large]
    \operatorname{Path}(\mathbb A^{\mathrm{op}})
    &
    \mathscr T_\Phi
        \arrow[l, "{\operatorname{Path}(I_\Phi)}"']
        \arrow[r, "{\operatorname{Path}(O_\Phi)}"]
    &
    \operatorname{Path}(\mathbb B^{\mathrm{op}}).
\end{tikzcd}
\]
\end{definition}

At path length \(n\), the left map records the input path, the right map
records the output-comparison path, and the middle object records the state
trajectory.

\subsection{Deterministic objects in the labelled topos}
\label{subsec:deterministic-objects-labelled-topos}

\begin{corollary}[Trajectory lifting pullback]
\label{cor:trajectory-lifting-pullback}
For every \(n\geq0\), the square
\[
\begin{tikzcd}
    \mathscr T_\Phi(n)
        \arrow[r, "{\operatorname{Path}(I_\Phi)_n}"]
        \arrow[d, "s_n^\Phi"']
    &
    \operatorname{Path}(\mathbb A^{\mathrm{op}})(n)
        \arrow[d, "{\lambda_0^{A,n}}"]
    \\
    P
        \arrow[r, "p_P"']
    &
    A_0
\end{tikzcd}
\]
is a pullback.
\end{corollary}

\begin{proof}
This is Proposition~\ref{prop:deterministic-path-lifting} under the
identification of finite input-path objects with the path objects of
\(\mathbb A^{\mathrm{op}}\).
\end{proof}

Thus a Mealy execution object is deterministic over its input label in every
finite degree. We use this as a structural property of a distinguished class
of objects in the labelled trajectory topos; no claim is made here that this
class is itself a subtopos.

\subsection{Composition of labelled finite trajectories}
\label{subsec:composition-trajectory-spans}

\begin{theorem}[Composition of labelled finite trajectories]
\label{thm:composition-labelled-finite-trajectories}
For composable Mealy morphisms
\[
    \Phi:\mathbb A\rightsquigarrow\mathbb B,
    \qquad
    \Psi:\mathbb B\rightsquigarrow\mathbb C,
\]
there is a canonical isomorphism
\[
    \mathscr T_{\Psi\circ\Phi}
    \cong
    \mathscr T_\Phi
    \times_{\operatorname{Path}(\mathbb B^{\mathrm{op}})}
    \mathscr T_\Psi.
\]
Equivalently, for every \(n\geq-1\),
\[
    \mathscr T_{\Psi\circ\Phi}(n)
    \cong
    \mathscr T_\Phi(n)
    \times_{\operatorname{Path}(\mathbb B^{\mathrm{op}})(n)}
    \mathscr T_\Psi(n).
\]
\end{theorem}

\begin{proof}
Apply the finite-limit-preserving graph path construction to
Theorem~\ref{thm:execution-pipeline-pullback}.
\end{proof}

Thus an \(n\)-step execution of the pipeline consists of:
\begin{enumerate}
    \item an \(n\)-step execution of \(\Phi\);
    \item an \(n\)-step execution of \(\Psi\);
    \item equality of the output-comparison path of \(\Phi\) and the input
    path of \(\Psi\).
\end{enumerate}
The equality is pathwise: the two chosen segmentations agree at every step.

\subsection{Strict and simulated trajectory maps}
\label{subsec:strict-simulated-trajectory-maps}

A strict semantic homomorphism
\[
    f:P\to Q
\]
induces a morphism
\[
    \mathscr T_f:
    \mathscr T_P\to\mathscr T_Q
\]
in the labelled trajectory topos.

A general Mealy simulation
\[
    (\theta,\alpha):P\Longrightarrow Q
\]
induces a trajectory morphism
\[
    \mathscr T_\theta:
    \mathscr T_Q\to\mathscr T_P
\]
that preserves input labels strictly. The output comparison is not equality
in the labelled slice. Instead, \(\alpha\) supplies a coherent comparison
between the two output-path maps. This comparison is promoted to a
simplicial natural transformation in
Section~\ref{subsec:simulations-simplicial-prisms}.

\section{Segal composition of finite paths}
\label{sec:segal-composition-finite-paths}

The trajectory sheaf of a graph remembers finite strings and contiguous
restriction. It does not by itself specify identities or composition. Those
operations are encoded by extending from contiguous injections to all
order-preserving maps.

\subsection{The augmented simplex category}
\label{subsec:augmented-simplex-category}

Let
\[
    \Delta_+
\]
be the augmented simplex category with objects
\[
    [-1],[0],[1],[2],\ldots
\]
and all order-preserving maps.

There is an identity-on-objects inclusion
\[
    j:
    \mathbf{Int}_{\mathrm{disc}}
    \hookrightarrow
    \Delta_+.
\]
Its nonempty morphisms are the distance-preserving injections.

The additional maps in \(\Delta_+\) encode category operations:
\begin{enumerate}
    \item codegeneracies insert identity arrows;
    \item inner cofaces contract adjacent arrows by composition.
\end{enumerate}
The simplicial identities encode unit and associativity.

\subsection{The augmented internal nerve}
\label{subsec:augmented-internal-nerve}

Let
\[
    \mathbb C=(C_0,C_1,s_C,t_C,e_C,m_C)
\]
be an internal category.

\begin{definition}[Augmented internal nerve]
\label{def:augmented-internal-nerve}
The \textbf{augmented nerve} is
\[
    N_+\mathbb C:
    \Delta_+^{\mathrm{op}}
    \to
    \mathcal E
\]
with
\[
    N_+\mathbb C_{-1}:=1,
    \qquad
    N_+\mathbb C_0:=C_0,
\]
and
\[
    N_+\mathbb C_n
    :=
    \underbrace{
    C_1\times_{C_0}\cdots\times_{C_0}C_1
    }_{n\text{ factors}}
\]
for \(n\geq1\).

Outer faces truncate paths, inner faces compose adjacent arrows, and
degeneracies insert identities. The unique map
\[
    [-1]\to[n]
\]
induces the unique augmentation
\[
    N_+\mathbb C_n\to N_+\mathbb C_{-1}=1.
\]
\end{definition}

\begin{definition}[Augmented strict Segal object]
\label{def:augmented-strict-segal-object}
An \textbf{augmented strict Segal object} is a functor
\[
    Y:\Delta_+^{\mathrm{op}}\to\mathcal E
\]
such that:
\begin{enumerate}
    \item \(Y_{-1}\to1\) is an isomorphism;
    \item for every \(n\geq2\), the Segal map
    \[
        Y_n
        \longrightarrow
        \underbrace{
        Y_1\times_{Y_0}\cdots\times_{Y_0}Y_1
        }_{n\text{ factors}}
    \]
    is an isomorphism.
\end{enumerate}
\end{definition}

Write
\[
    \mathsf{Seg}^{\mathrm{str}}_+(\mathcal E)
\]
for the full subcategory of augmented strict Segal objects.

\begin{theorem}[Internal nerve equivalence]
\label{thm:internal-nerve-equivalence}
The augmented nerve defines an equivalence
\[
\begin{tikzcd}[column sep=large]
    \mathsf{IntCat}(\mathcal E)
        \arrow[r, shift left=0.7ex, "N_+"]
    &
    \mathsf{Seg}^{\mathrm{str}}_+(\mathcal E)
        \arrow[l, shift left=0.7ex, "R"] .
\end{tikzcd}
\]
\end{theorem}

\begin{proof}
Let \(Y\) be augmented strict Segal. Define
\[
    R(Y)_0:=Y_0,
    \qquad
    R(Y)_1:=Y_1.
\]
The two endpoint inclusions define source and target. The unique codegeneracy
\[
    [1]\to[0]
\]
defines the identity map
\[
    Y_0\to Y_1.
\]
The inner coface
\[
    [1]\to[2],
    \qquad
    0\mapsto0,
    \quad
    1\mapsto2,
\]
defines composition after transport across
\[
    Y_2
    \cong
    Y_1\times_{Y_0}Y_1.
\]
The simplicial identities give the internal-category equations.

Conversely, the nerve of an internal category is strict Segal by its
iterated-pullback description. Reconstruction recovers the category in
degrees zero and one, while the Segal isomorphisms identify all higher
degrees.

A natural transformation between strict Segal nerves is determined by its
degree-zero and degree-one components. Compatibility with the endpoint,
degeneracy, and inner-face maps says exactly that these components define an
internal functor. Hence the nerve is fully faithful as well as essentially
surjective.
\end{proof}

\subsection{Underlying trajectory sheaves}
\label{subsec:underlying-trajectory-sheaves}

\begin{definition}[Underlying trajectory sheaf]
\label{def:underlying-trajectory-sheaf}
The \textbf{underlying trajectory sheaf} of an internal category
\(\mathbb C\) is
\[
    \operatorname{Path}(\mathbb C)
    :=
    j^\ast N_+\mathbb C.
\]
\end{definition}

If \(U\mathbb C\) is the underlying internal graph, then
\[
    \operatorname{Path}(\mathbb C)
    \cong
    \operatorname{Path}(U\mathbb C),
\]
since both sides consist of finite composable strings with contiguous
subpath restriction.

The restriction \(j^\ast N_+\mathbb C\) remembers:
\begin{enumerate}
    \item objects and arrows;
    \item finite composable strings;
    \item contiguous restriction;
    \item gluing along matching boundaries.
\end{enumerate}
The full augmented nerve additionally remembers:
\begin{enumerate}
    \item identity insertion;
    \item contraction by composition;
    \item unit and associativity coherence.
\end{enumerate}

\subsection{Segal composition structures}
\label{subsec:segal-composition-structures}

Let \(X\) be a trajectory sheaf. Write
\[
    s_X:=X(\epsilon_0),
    \qquad
    t_X:=X(\epsilon_1).
\]
For \(n\geq1\), let
\[
    \operatorname{seg}_{X,n}:
    X(n)
    \xrightarrow{\cong}
    \underbrace{
    X(1)\times_{X(0)}\cdots\times_{X(0)}X(1)
    }_{n\text{ factors}}
\]
be the elementary-edge restriction isomorphism, with inverse
\[
    \operatorname{gl}_{X,n}.
\]
Set
\[
    \operatorname{seg}_{X,0}
    =
    \operatorname{gl}_{X,0}
    =
    1_{X(0)}.
\]
Write
\[
    \rho_1,\rho_2:X(2)\rightrightarrows X(1)
\]
for the two elementary-edge restrictions, and
\[
    \nu_i^{X,n}:X(n)\to X(0)
\]
for the vertex restrictions.

\begin{definition}[Segal composition structure]
\label{def:segal-composition-structure}
A \textbf{Segal composition structure} on \(X\) consists of maps
\[
    e_X:X(0)\to X(1),
    \qquad
    c_X:X(2)\to X(1)
\]
satisfying the following equations.

First, identities have the correct endpoints:
\[
    s_Xe_X=1_{X(0)},
    \qquad
    t_Xe_X=1_{X(0)}.
\]

Second, contraction has the endpoints of the original two-edge path:
\[
    s_Xc_X=\nu_0^{X,2},
    \qquad
    t_Xc_X=\nu_2^{X,2}.
\]

Define identity-insertion maps
\[
    \ell_X
    :=
    \operatorname{gl}_{X,2}
    \left\langle
        e_Xs_X,
        1_{X(1)}
    \right\rangle,
\]
\[
    r_X
    :=
    \operatorname{gl}_{X,2}
    \left\langle
        1_{X(1)},
        e_Xt_X
    \right\rangle.
\]
The unit equations are
\[
    c_X\ell_X=1_{X(1)},
    \qquad
    c_Xr_X=1_{X(1)}.
\]

Finally, under
\[
    X(3)
    \cong
    X(1)\times_{X(0)}X(1)\times_{X(0)}X(1),
\]
the two morphisms \(X(3)\to X(1)\) obtained by first contracting the
first two edges and by first contracting the last two edges are equal.
\end{definition}

Equivalently, under
\[
    X(2)
    \cong
    X(1)\times_{X(0)}X(1),
\]
the map \(c_X\) defines an internal-category composition on the graph
\(G_X\), with identity \(e_X\).

\begin{theorem}[Classification of Segal composition structures]
\label{thm:classification-segal-composition-structures}
For a trajectory sheaf \(X\), there are canonical bijections between:
\begin{enumerate}
    \item Segal composition structures on \(X\);
    \item internal-category structures on the reconstructed graph \(G_X\);
    \item isomorphism classes over \(X\) of pairs \((Y,\vartheta)\), where
    \(Y\) is an augmented strict Segal object and
    \[
        \vartheta:j^\ast Y\xrightarrow{\cong}X
    \]
    is a specified isomorphism.
\end{enumerate}
\end{theorem}

\begin{proof}
The graph--trajectory-sheaf equivalence identifies the degree-zero and
degree-one objects of \(X\) with the vertices and edges of \(G_X\). Under
the degree-two sheaf isomorphism, the displayed identity, endpoint, unit, and
associativity equations are precisely the internal-category equations.

The internal nerve equivalence identifies internal-category structures on
\(G_X\) with augmented strict Segal objects whose inert restriction is the
path sheaf of \(G_X\), hence with augmented strict Segal objects equipped
with an identification of their restriction with \(X\).
\end{proof}

\begin{corollary}[Canonical Segal composition of execution]
\label{cor:canonical-segal-composition-execution}
Every internal category, and in particular every
\(\operatorname{Exec}(\Phi)\), determines:
\begin{enumerate}
    \item an underlying trajectory sheaf;
    \item a canonical Segal composition structure supplied by its augmented
    nerve.
\end{enumerate}
\end{corollary}

\subsection{Identity insertion, contraction, and segmentation}
\label{subsec:execution-identity-contraction}

At a temporal boundary \(x_i\), the degeneracy map inserts the identity
execution
\[
\begin{tikzcd}[column sep=huge]
    x_i
        \arrow[r, "{(e_A(p_Px_i),x_i)}"]
    &
    x_i.
\end{tikzcd}
\]

A two-step execution
\[
\begin{tikzcd}[column sep=large]
    x
        \arrow[r, "{(b,x)}"]
    &
    \delta_P(b,x)
        \arrow[r, "{(a,\delta_P(b,x))}"]
    &
    \delta_P\bigl(a,\delta_P(b,x)\bigr)
\end{tikzcd}
\]
contracts to
\[
\begin{tikzcd}[column sep=huge]
    x
        \arrow[r, "{(b\circ a,x)}"]
    &
    \delta_P(b\circ a,x).
\end{tikzcd}
\]
The transition multiplication law identifies the terminal states, while the
output multiplication law identifies the output label of the contracted step
with the composite of the original output labels.

For \(n\geq1\), total contraction is the iterated composition map
\[
    \kappa_n^\Phi:
    \mathscr T_\Phi(n)
    \longrightarrow
    \mathscr T_\Phi(1).
\]
For \(n=0\), put
\[
    \kappa_0^\Phi:=e_{\operatorname{Exec}}:P\to D_P.
\]
These maps satisfy
\[
    s_1^\Phi\kappa_n^\Phi=s_n^\Phi,
    \qquad
    t_1^\Phi\kappa_n^\Phi=t_n^\Phi.
\]
For \(0\leq i\leq n\), prefix contraction similarly gives
\[
    \kappa_{n,i}^\Phi:
    \mathscr T_\Phi(n)
    \longrightarrow
    \mathscr T_\Phi(1)
\]
with
\[
    s_1^\Phi\kappa_{n,i}^\Phi=s_n^\Phi,
    \qquad
    t_1^\Phi\kappa_{n,i}^\Phi=\nu_i^{\Phi,n}.
\]
For \(i=0\), this inserts the identity at the initial state.

\begin{remark}[Path length records segmentation]
\label{rem:path-length-segmentation}
The integer \(n\) records a chosen decomposition of an execution into
\(n\) categorical arrows. It does not generally measure irreducible
duration or a number of primitive computational steps.

Every composable string
\[
    x_0\to x_1\to\cdots\to x_n
\]
contracts to a one-step execution with the same initial and terminal states.
The finite-path objects retain the chosen segmentation, while contraction
records the corresponding categorical composite.

Consequently, unconstrained endpoint reachability and universal endpoint
safety for the full categorical execution system are already detected by
one-step executions: a one-step input label may itself be an arbitrary
composite arrow of \(\mathbb A\). A primitive-step interpretation requires
a chosen generating graph or another specified class of atomic transitions.
\end{remark}

\subsection{Simulations as simplicial prisms}
\label{subsec:simulations-simplicial-prisms}

Let
\[
    (\theta,\alpha):P\Longrightarrow Q
\]
be a Mealy simulation. The execution functor induces a simplicial map
\[
    N_+\operatorname{Exec}(\theta):
    N_+\operatorname{Exec}(Q)
    \longrightarrow
    N_+\operatorname{Exec}(P).
\]
The natural transformation
\[
    \alpha^{\mathrm{op}}:
    O_Q\Longrightarrow O_P\operatorname{Exec}(\theta)
\]
determines the standard simplicial prism between
\[
    N_+O_Q
\]
and
\[
    N_+O_P\,N_+\operatorname{Exec}(\theta).
\]
In degree \(n\), this prism compares the two output paths by the boundary
components
\[
    \alpha_{x_0}^{\mathrm{op}},
    \alpha_{x_1}^{\mathrm{op}},
    \ldots,
    \alpha_{x_n}^{\mathrm{op}},
\]
with each square commuting by the simulation equation.

Thus simulations preserve the entire segmented input execution strictly and
compare the output-labelled simplices coherently at every boundary.

\subsection{Segal-enhanced labelled execution}
\label{subsec:segal-enhanced-labelled-execution}

\begin{definition}[Segal-enhanced labelled execution]
\label{def:segal-enhanced-labelled-execution}
The \textbf{Segal-enhanced labelled execution semantics} of \(\Phi\) is the
span of augmented strict Segal objects
\[
\begin{tikzcd}[column sep=large]
    N_+(\mathbb A^{\mathrm{op}})
    &
    N_+\operatorname{Exec}(\Phi)
        \arrow[l, "N_+I_\Phi"']
        \arrow[r, "N_+O_\Phi"]
    &
    N_+(\mathbb B^{\mathrm{op}}).
\end{tikzcd}
\]
\end{definition}

Restriction along
\[
    j:\mathbf{Int}_{\mathrm{disc}}\hookrightarrow\Delta_+
\]
recovers the labelled finite-trajectory object. The trajectory object records
finite paths, restriction, and gluing. The Segal enhancement additionally
records identity insertion, contraction, and coherence of iterated
contraction.

\subsection{Finite-limit preservation}
\label{subsec:path-nerve-finite-limits}

\begin{proposition}[Finite-limit preservation]
\label{prop:path-nerve-finite-limit-preservation}
The functors
\[
    N_+:
    \mathsf{IntCat}(\mathcal E)
    \to
    [\Delta_+^{\mathrm{op}},\mathcal E]
\]
and
\[
    \operatorname{Path}:
    \mathsf{IntCat}(\mathcal E)
    \to
    \mathbf{TrajSh}_{\mathcal E}
\]
preserve finite limits.
\end{proposition}

\begin{proof}
Finite limits of internal categories are computed on object and arrow
objects. Every nerve degree is obtained from those objects by finite iterated
pullbacks, so \(N_+\) preserves finite limits degreewise. Restriction along
\(j\) preserves pointwise limits, and limits of trajectory sheaves are
computed pointwise.
\end{proof}

\section{Decomposition-space execution semantics}
\label{sec:decomposition-space-execution}

The Segal structure records how a string of execution arrows contracts to a
composite. The same nerve also records the complementary operation: how a
single execution arrow may be factored into a string of composable
executions. This is the decomposition-space aspect of execution semantics.

\subsection{Active and inert simplicial maps}
\label{subsec:active-inert-maps}

Restrict the augmented nerve along the inclusion
\[
    \Delta\hookrightarrow\Delta_+
\]
and write
\[
    N\mathbb C:\Delta^{\mathrm{op}}\to\mathcal E
\]
for the ordinary nerve.

A morphism of \(\Delta\) is \textbf{inert} when it is distance-preserving;
it selects a contiguous subinterval. A morphism is \textbf{active} when it
preserves the two endpoints. Every simplicial map factors uniquely as an
active map followed by an inert map.

Inert maps therefore encode temporal restriction, while active maps encode
contraction and factorization of the interval as a whole.

\begin{definition}[Internal decomposition object]
\label{def:internal-decomposition-object}
A simplicial object
\[
    X:\Delta^{\mathrm{op}}\to\mathcal E
\]
is an \textbf{internal decomposition object} if it sends every active--inert
pushout square in \(\Delta\) to a pullback square in \(\mathcal E\).
\end{definition}

\subsection{Execution nerves are decomposition objects}
\label{subsec:execution-nerve-decomposition}

\begin{theorem}[Segal objects are decomposition objects]
\label{thm:segal-implies-decomposition}
Every strict Segal object in \(\mathcal E\) is an internal decomposition
object. In particular, for every internal category \(\mathbb C\),
\[
    N\mathbb C
\]
is an internal decomposition object, and hence so is
\[
    N\operatorname{Exec}(\Phi).
\]
\end{theorem}

\begin{proof}
For a strict Segal object, every degree is an iterated pullback of the
one-simplex object over the zero-simplex object. Under these identifications,
the image of an active--inert pushout square is a pullback square selecting
and composing the corresponding consecutive blocks. Thus the decomposition
pullback conditions follow from the stronger Segal pullback identities.
\end{proof}

The Segal and decomposition readings are complementary:
\begin{enumerate}
    \item the Segal reading asks for the composite represented by a segmented
    execution;
    \item the decomposition reading asks for the coherent space or object of
    segmentations of a given execution.
\end{enumerate}

\subsection{The objective incidence comonoid}
\label{subsec:objective-incidence-comonoid}

Let \(X\) be a decomposition object. Its objective incidence
comultiplication is represented by the span
\[
\begin{tikzcd}[column sep=large]
    X_1
    &
    X_2
        \arrow[l,"d_1"']
        \arrow[r,"{(d_2,d_0)}"]
    &
    X_1\times X_1.
\end{tikzcd}
\]
Here \(d_1\) contracts a two-step path, while \((d_2,d_0)\) remembers its
two factors. The counit is represented by
\[
\begin{tikzcd}[column sep=large]
    X_1
    &
    X_0
        \arrow[l,"s_0"']
        \arrow[r]
    &
    1.
\end{tikzcd}
\]

\begin{theorem}[Execution incidence comonoid]
\label{thm:execution-incidence-comonoid}
For every internal category \(\mathbb C\), the spans
\[
    C_1
    \xleftarrow{\ d_1\ }
    N\mathbb C_2
    \xrightarrow{\ (d_2,d_0)\ }
    C_1\times C_1
\]
and
\[
    C_1
    \xleftarrow{\ s_0\ }
    C_0
    \longrightarrow
    1
\]
define a comonoid object on \(C_1\) in the bicategory
\(\operatorname{Span}(\mathcal E)\).
\end{theorem}

\begin{proof}
The two iterated comultiplications classify the two ways of cutting an arrow
into three consecutive factors. The decomposition pullback conditions
identify both with \(N\mathbb C_3\). Hence coassociativity follows from the
simplicial identities. The counit equations follow from the degeneracy and
outer-face identities.
\end{proof}

For
\[
    \mathbb C=\operatorname{Exec}(\Phi),
\]
the comultiplication sends an execution arrow to the object of all its binary
execution factorizations. Unlike total contraction, this retains the chosen
cut and the intermediate state.

\subsection{Unique lifting of execution factorizations}
\label{subsec:unique-lifting-factorizations}

A functor \(F:\mathbb C\to\mathbb D\) is said to have
\textbf{unique lifting of factorizations} if every factorization of
\(F(f)\) in \(\mathbb D\) lifts uniquely to a factorization of \(f\) in
\(\mathbb C\). It is \textbf{conservative} if it reflects identity arrows.
A conservative functor with unique lifting of factorizations is called
\textbf{CULF}.

\begin{proposition}[Simulation functors are CULF]
\label{prop:simulation-functors-culf}
For every Mealy simulation
\[
    (\theta,\alpha):P\Longrightarrow Q,
\]
the induced functor
\[
    \operatorname{Exec}(\theta):
    \operatorname{Exec}(Q)
    \longrightarrow
    \operatorname{Exec}(P)
\]
is CULF.
\end{proposition}

\begin{proof}
The functor preserves the input arrow exactly. If
\[
    \operatorname{Exec}(\theta)(a,y)
\]
is an identity, then its input label \(a\) is an identity, and hence
\((a,y)\) is an identity. Thus the functor is conservative.

A factorization of
\[
    (a,\theta y)
\]
in \(\operatorname{Exec}(P)\) determines a factorization of the input
arrow \(a^{\mathrm{op}}\) in \(\mathbb A^{\mathrm{op}}\). Since
\[
    I_Q:\operatorname{Exec}(Q)\to\mathbb A^{\mathrm{op}}
\]
is a discrete opfibration, that input factorization lifts uniquely from the
source state \(y\). Transition equivariance ensures that its image under
\(\operatorname{Exec}(\theta)\) is the prescribed factorization in
\(\operatorname{Exec}(P)\). Hence factorizations lift uniquely.
\end{proof}

\begin{corollary}[Simulation morphisms of incidence comonoids]
\label{cor:simulation-incidence-comonoid-map}
Every Mealy simulation induces, in its displayed direction, a morphism of
objective execution incidence comonoids.
\end{corollary}

\begin{proof}
The nerve of a CULF functor is cartesian on the active maps controlling
factorization. Therefore it intertwines the incidence comultiplication and
counit spans.
\end{proof}

\subsection{Locally finite linearization}
\label{subsec:locally-finite-linearization}

The objective incidence comonoid exists without finiteness assumptions. If
\(\mathcal E=\mathbf{Set}\) and each execution arrow has finitely many
factorizations, one may apply free-module or cardinality constructions to
obtain an ordinary incidence coalgebra. Additional finite-length hypotheses
may support Möbius inversion. These finiteness conditions are not part of
the standing theory of this chapter.

\section{Trajectory quotients and behavioural minimization}
\label{sec:trajectory-quotients-minimization}

\subsection{Levelwise preservation of regular quotients}
\label{subsec:levelwise-preservation-regular-quotients}

Let
\[
    e:P\twoheadrightarrow Q
\]
be a regular-epimorphic strict semantic homomorphism between Mealy morphisms
over the same interfaces.

\begin{theorem}[Trajectory preservation of regular quotients]
\label{thm:trajectory-preservation-regular-quotients}
For every \(n\geq-1\),
\[
    \mathscr T_e(n):
    \mathscr T_P(n)
    \twoheadrightarrow
    \mathscr T_Q(n)
\]
is a regular epimorphism.
\end{theorem}

\begin{proof}
At degree \(-1\), the map is the identity of \(1\). For \(n\geq0\),
deterministic lifting gives
\[
    \mathscr T_P(n)
    \cong
    P\times_{A_0}
    \operatorname{Path}(\mathbb A^{\mathrm{op}})(n),
\]
\[
    \mathscr T_Q(n)
    \cong
    Q\times_{A_0}
    \operatorname{Path}(\mathbb A^{\mathrm{op}})(n).
\]
Under these identifications,
\[
    \mathscr T_e(n)
    =
    e\times_{A_0}1.
\]
This is a pullback of \(e\), hence a regular epimorphism.
\end{proof}

\begin{corollary}[Levelwise quotient of trajectory sheaves]
\label{cor:levelwise-quotient-trajectory-sheaves}
A regular-epimorphic strict semantic homomorphism induces a levelwise regular
epimorphism of trajectory sheaves.
\end{corollary}

No stronger claim about regular epimorphisms in the entire trajectory-sheaf
category is required here.

\subsection{Compatibility with Segal structure}
\label{subsec:quotient-segal-compatibility}

\begin{proposition}[Segal compatibility of strict quotients]
\label{prop:strict-quotient-segal-compatibility}
The simplicial map
\[
    N_+\operatorname{Exec}(e):
    N_+\operatorname{Exec}(P)
    \longrightarrow
    N_+\operatorname{Exec}(Q)
\]
commutes with all restrictions, degeneracies, and face maps. Consequently,
the trajectory quotient preserves contiguous restriction, gluing, identity
insertion, and contraction.
\end{proposition}

\begin{proof}
The map is the nerve of an internal functor. Naturality with respect to every
map of \(\Delta_+\) gives the result.
\end{proof}

\subsection{Compatibility with factorization}
\label{subsec:quotient-factorization-compatibility}

Regarded in the established simulation orientation, the strict homomorphism
\[
    e:P\to Q
\]
defines a strict simulation
\[
    Q\Longrightarrow P.
\]
Hence
\[
    \operatorname{Exec}(e):
    \operatorname{Exec}(P)
    \to
    \operatorname{Exec}(Q)
\]
is CULF by
Proposition~\ref{prop:simulation-functors-culf}.

\begin{corollary}[Quotients preserve execution factorization]
\label{cor:quotients-preserve-factorization}
A regular-epimorphic strict semantic homomorphism induces a morphism of
objective execution incidence comonoids
\[
    \operatorname{IncExec}(P)
    \longrightarrow
    \operatorname{IncExec}(Q).
\]
\end{corollary}

Thus a semantic quotient preserves not only the image of every finite
trajectory but also the coherent factorization structure of execution
arrows.

\subsection{Behavioural trajectory minimization}
\label{subsec:behavioural-trajectory-minimization}

Recall the semantic-image factorization
\[
    P
    \xrightarrow{\eta_P}
    \operatorname{Beh}(P)
    \hookrightarrow
    \mathbf{Beh}^{\mathrm{can}}_{\mathbb A,\mathbb B},
\]
where
\[
    \eta_P:P\twoheadrightarrow\operatorname{Beh}(P)
\]
is a regular-epimorphic strict semantic homomorphism.

\begin{theorem}[Behavioural quotient of finite trajectories]
\label{thm:behavioural-quotient-finite-trajectories}
Every Mealy morphism \(P\) admits a canonical levelwise regular quotient
\[
    \mathscr T_P
    \twoheadrightarrow
    \mathscr T_{\operatorname{Beh}(P)}.
\]
At degree \(n\geq0\), its kernel pair is canonically
\[
    \operatorname{Eq}(\eta_P)
    \times_{A_0}
    \operatorname{Path}(\mathbb A^{\mathrm{op}})(n).
\]
Thus, internally, two executions have the same image precisely when they
have the same input path and their initial states have the same residual
semantic image. Equivalently, along a common input path, their semantic state
images agree at every boundary.
\end{theorem}

\begin{proof}
Apply
Theorem~\ref{thm:trajectory-preservation-regular-quotients}
to \(\eta_P\).

Under deterministic lifting, the degree-\(n\) map is
\[
    \eta_P\times_{A_0}1.
\]
Its kernel pair is therefore the displayed pullback of
\(\operatorname{Eq}(\eta_P)\), and strict transition preservation forces
equality of semantic images at every subsequent boundary.
\end{proof}

\begin{corollary}[Separated trajectory semantics]
\label{cor:separated-trajectory-semantics}
If \(P\) is behaviourally separated, then
\[
    \mathscr T_P
    \cong
    \mathscr T_{\operatorname{Beh}(P)}.
\]
\end{corollary}

\begin{proof}
For a behaviourally separated system, the semantic quotient
\(\eta_P\) is an isomorphism.
\end{proof}

Behavioural separatedness removes duplicate trajectories arising from
behaviourally indistinguishable initial states. It does not identify
different segmentations; those are related by Segal contraction and retained
as distinct factorization data by the incidence comonoid.

\begin{theorem}[Behavioural minimization of structured execution]
\label{thm:behavioural-minimization-structured-execution}
The semantic quotient
\[
    P\twoheadrightarrow\operatorname{Beh}(P)
\]
induces simultaneously:
\begin{enumerate}
    \item a levelwise regular quotient of finite trajectory sheaves;
    \item a morphism of augmented strict Segal execution objects;
    \item a morphism of objective execution incidence comonoids.
\end{enumerate}
\end{theorem}

\begin{proof}
The three assertions are respectively
Theorem~\ref{thm:trajectory-preservation-regular-quotients},
Proposition~\ref{prop:strict-quotient-segal-compatibility}, and
Corollary~\ref{cor:quotients-preserve-factorization}.
\end{proof}

\subsection{Finite-trajectory Myhill--Nerode theory}
\label{subsec:finite-trajectory-myhill-nerode}

Let
\[
    k:K\to
    \mathbf{Beh}_{\mathbb A,\mathbb B,0}^{\mathrm{can}}
\]
be a behaviour specification, and let
\[
    \operatorname{Min}(K)
\]
be its minimal residual realization, namely the restriction of the canonical
behaviour machine to
\[
    \operatorname{Der}_0(K).
\]

\begin{corollary}[Finite-trajectory Myhill--Nerode theorem]
\label{cor:finite-trajectory-myhill-nerode}
Let
\[
    i:K\to P
\]
be a reachable realization. The canonical semantic quotient
\[
    P\twoheadrightarrow\operatorname{Min}(K)
\]
induces:
\begin{enumerate}
    \item a levelwise regular epimorphism
    \[
        \mathscr T_P
        \twoheadrightarrow
        \mathscr T_{\operatorname{Min}(K)};
    \]
    \item a morphism of augmented strict Segal execution objects;
    \item a morphism of objective execution incidence comonoids.
\end{enumerate}
If \(P\) is also behaviourally separated, the trajectory and Segal maps are
isomorphisms.
\end{corollary}

\begin{proof}
Apply
Theorem~\ref{thm:behavioural-minimization-structured-execution}
to the semantic quotient supplied by the categorical Myhill--Nerode theorem.
In the separated case, that quotient is an isomorphism.
\end{proof}

Thus reachable and separated realizations have exactly the labelled finite
trajectories and Segal composition structure of the canonical derivative
realization. Their execution factorization comonoids are correspondingly
isomorphic.

\section{The temporal modal calculus of trajectories}
\label{sec:temporal-modal-calculus}

The Grothendieck topos of trajectory sheaves carries an internal
higher-order logic. The source and target maps of finite trajectory objects
therefore induce canonical temporal modalities on state predicates.

\subsection{The trajectory subobject hyperdoctrine}
\label{subsec:trajectory-subobject-hyperdoctrine}

Let \(X\) be a trajectory sheaf. For every \(n\geq0\), write
\[
    s_n,t_n:X(n)\rightrightarrows X(0)
\]
for its initial and terminal boundary maps. Since
\(\mathbf{TrajSh}_{\mathcal E}\) is a topos, pullback on subobjects has both
adjoints:
\[
    \exists_{s_n}
    \dashv
    s_n^\ast
    \dashv
    \forall_{s_n},
\]
and similarly for \(t_n\).

All subobjects and adjunctions in this section may equivalently be computed
in \(\mathcal E\) degreewise. The topos formulation records their stability
under reindexing and their compatibility with the full labelled-trajectory
logic.

\subsection{Endpoint temporal modalities}
\label{subsec:endpoint-temporal-modalities}

Let
\[
    C\hookrightarrow X(0)
\]
be a state predicate.

\begin{definition}[Existential future]
\label{def:existential-future}
The \textbf{existential \(n\)-segment future} of \(C\) is
\[
    \Diamond_n C
    :=
    \exists_{s_n}(t_n^\ast C).
\]
\end{definition}

Internally,
\[
    x\in\Diamond_nC
\]
if and only if there exists an \(n\)-segment trajectory beginning at \(x\)
and ending in \(C\).

\begin{definition}[Universal future]
\label{def:universal-future}
The \textbf{universal \(n\)-segment future} of \(C\) is
\[
    \Box_n C
    :=
    \forall_{s_n}(t_n^\ast C).
\]
\end{definition}

Internally,
\[
    x\in\Box_nC
\]
if and only if every \(n\)-segment trajectory beginning at \(x\) ends in
\(C\).

\begin{definition}[Forward endpoint image]
\label{def:forward-endpoint-image}
The \textbf{forward \(n\)-segment endpoint image} of \(C\) is
\[
    \operatorname{Post}_n(C)
    :=
    \exists_{t_n}(s_n^\ast C).
\]
\end{definition}

Internally,
\[
    y\in\operatorname{Post}_n(C)
\]
if and only if there exists an \(n\)-segment trajectory beginning in \(C\)
and ending at \(y\).

The operators \(\Diamond_n\) and \(\operatorname{Post}_n\) preserve joins,
while \(\Box_n\) preserves meets. At length zero,
\[
    \Diamond_0
    =
    \Box_0
    =
    \operatorname{Post}_0
    =
    1.
\]

\subsection{Boundary-sensitive safety}
\label{subsec:boundary-sensitive-safety}

For \(0\leq i\leq n\), define
\[
    \Box_{n,i}C
    :=
    \forall_{s_n}
    \left(
        (\nu_i^{X,n})^\ast C
    \right).
\]
Thus \(\Box_{n,i}C\) consists of states from which every \(n\)-segment
trajectory lies in \(C\) at boundary \(i\).

Define the finite-prefix safety predicate
\[
    \Box_{\mathrm{pref}}C
    :=
    \bigwedge_{n\in\mathbb N}
    \bigwedge_{0\leq i\leq n}
    \Box_{n,i}C.
\]
This is the largest state predicate from which all finite trajectories remain
in \(C\) at every displayed boundary.

\subsection{Label-constrained modalities}
\label{subsec:label-constrained-modalities}

Suppose \(X\) is labelled over
\(\operatorname{Path}(\mathbb A^{\mathrm{op}})\), with degree-\(n\) input
label
\[
    L_n:X(n)\to\operatorname{Path}(\mathbb A^{\mathrm{op}})(n).
\]
Let
\[
    U\hookrightarrow
    \operatorname{Path}(\mathbb A^{\mathrm{op}})(n)
\]
be an input-path predicate.

\begin{definition}[Action-indexed temporal modalities]
\label{def:action-indexed-temporal-modalities}
Define
\[
    \langle U\rangle C
    :=
    \exists_{s_n}
    \left(
        L_n^\ast U
        \wedge
        t_n^\ast C
    \right),
\]
and
\[
    [U]C
    :=
    \forall_{s_n}
    \left(
        L_n^\ast U
        \Rightarrow
        t_n^\ast C
    \right).
\]
\end{definition}

Internally, \(\langle U\rangle C\) asserts the existence of a trajectory
whose input path satisfies \(U\) and whose terminal state satisfies \(C\).
The predicate \([U]C\) asserts that every trajectory whose input path
satisfies \(U\) ends in \(C\).

Output constraints are incorporated in the same way using the output-label
map. Since the labelled trajectory category is a slice topos, input, output,
and state predicates may be combined by the full higher-order internal
logic.

\subsection{Modal composition from temporal gluing}
\label{subsec:modal-composition-gluing}

\begin{theorem}[Composition laws for temporal modalities]
\label{thm:composition-laws-temporal-modalities}
For every trajectory sheaf \(X\) and \(m,n\in\mathbb N\), temporal gluing
and Beck--Chevalley give
\[
    \Diamond_{m+n}
    =
    \Diamond_m\Diamond_n,
\]
\[
    \Box_{m+n}
    =
    \Box_m\Box_n,
\]
and
\[
    \operatorname{Post}_{m+n}
    =
    \operatorname{Post}_n\operatorname{Post}_m.
\]
\end{theorem}

\begin{proof}
Use the canonical pullback
\[
    X(m+n)
    \cong
    X(m)\times_{X(0)}X(n),
\]
where the terminal boundary of the first path is identified with the initial
boundary of the second. Dependent sum composition and Beck--Chevalley give
the two existential formulas. Dependent product composition and
Beck--Chevalley give the universal formula. The displayed order reflects that
an \(m\)-segment path is executed before an \(n\)-segment path.
\end{proof}

\subsection{Categorical collapse and primitive time}
\label{subsec:categorical-collapse-primitive-time}

Let \(X=\operatorname{Path}(\mathbb C)\) arise from an internal category.
Total contraction gives a map
\[
    \kappa_n:X(n)\to X(1)
\]
with the same endpoints. Conversely, every one-step arrow expands to an
\(n\)-step path by inserting identities.

\begin{proposition}[Collapse of unconstrained endpoint modalities]
\label{prop:collapse-unconstrained-modalities}
For every \(n\geq1\),
\[
    \Diamond_n=\Diamond_1,
    \qquad
    \Box_n=\Box_1,
    \qquad
    \operatorname{Post}_n=\operatorname{Post}_1.
\]
\end{proposition}

\begin{proof}
Contraction shows that every \(n\)-step endpoint pair is a one-step endpoint
pair. Identity insertion shows that every one-step endpoint pair is realized
by an \(n\)-step path. Hence the corresponding endpoint relations, and
therefore their existential and universal predicate transformers, agree.
\end{proof}

This collapse applies only to unconstrained endpoint modalities. The
following remain nontrivial:
\begin{enumerate}
    \item exact input-path constraints;
    \item output-path constraints;
    \item segmentation-sensitive predicates;
    \item intermediate-boundary formulas;
    \item modalities on a chosen primitive generating graph.
\end{enumerate}

\section{Residual execution and the derivative modal calculus}
\label{sec:residual-execution-derivative-modal-calculus}

We now specialize the trajectory modalities to the canonical behaviour
machine. This identifies the generated residual closure and residual safety
core with the forward and universal temporal operators.

\subsection{Canonical residual trajectories}
\label{subsec:canonical-residual-trajectories}

Fix internal categories
\[
    \mathbb A,
    \qquad
    \mathbb B,
\]
and let
\[
    \mathbf{Beh}^{\mathrm{can}}_{\mathbb A,\mathbb B}
\]
be the canonical behaviour Mealy morphism, with carrier
\[
    \mathbf{Beh}_0.
\]
Write
\[
    \mathscr T_{\mathbf{Beh}}
    :=
    \mathscr T_{
        \mathbf{Beh}^{\mathrm{can}}_{\mathbb A,\mathbb B}
    }.
\]

For \(n\geq0\) and \(0\leq i\leq n\), write
\[
    \nu_i^{(n)}:
    \mathscr T_{\mathbf{Beh}}(n)
    \to
    \mathbf{Beh}_0
\]
for the \(i\)-th residual-state map. Put
\[
    s_n:=\nu_0^{(n)},
    \qquad
    t_n:=\nu_n^{(n)}.
\]

A generalized residual execution has states
\[
    H_0,H_1,\ldots,H_n
\]
with
\[
    H_i=\partial_{a_i}H_{i-1}.
\]

The one-step trajectory object is canonically the derivative domain:
\[
    \mathscr T_{\mathbf{Beh}}(1)
    \cong
    D\mathbf{Beh}
    =
    A_1\times_{t_A,A_0,p_{\mathbf{Beh}}}\mathbf{Beh}_0.
\]
Under this identification,
\[
    s_1=\operatorname{pr}_\partial,
    \qquad
    t_1=\partial.
\]
Pullbacks, images, and dependent products below are computed in
\(\mathcal E\), since initial and derivative states may have different
anchor coordinates.

\subsection{Trajectories of an invariant predicate}
\label{subsec:trajectories-invariant-predicate}

Let
\[
    V\hookrightarrow\mathbf{Beh}_0
\]
be residually invariant, equivalently derivative-closed, and let
\(\iota(V)\) denote the restricted canonical Mealy realization.

\begin{proposition}[Trajectories of an invariant predicate]
\label{prop:trajectories-invariant-predicate}
For every \(n\geq0\), there is a canonical pullback isomorphism
\[
    \mathscr T_{\iota(V)}(n)
    \cong
    V\times_{\mathbf{Beh}_0,s_n}
    \mathscr T_{\mathbf{Beh}}(n).
\]
Consequently, a canonical residual trajectory belongs to
\(\mathscr T_{\iota(V)}(n)\) if and only if its initial state belongs to
\(V\).
\end{proposition}

\begin{proof}
Deterministic lifting gives
\[
    \mathscr T_{\mathbf{Beh}}(n)
    \cong
    \mathbf{Beh}_0
    \times_{A_0}
    \operatorname{Path}(\mathbb A^{\mathrm{op}})(n),
\]
and
\[
    \mathscr T_{\iota(V)}(n)
    \cong
    V
    \times_{A_0}
    \operatorname{Path}(\mathbb A^{\mathrm{op}})(n).
\]
The second is the pullback of the first along
\(V\hookrightarrow\mathbf{Beh}_0\). Residual invariance guarantees that
all later states remain in \(V\).
\end{proof}

\subsection{Residual invariance as modal fixed point}
\label{subsec:residual-invariance-modal-fixed-point}

\begin{theorem}[Finite-path characterization of residual invariance]
\label{thm:finite-path-characterization-derivative-invariance}
For a behavioural predicate
\[
    C\hookrightarrow\mathbf{Beh}_0,
\]
the following are equivalent:
\begin{enumerate}
    \item \(C\) is residually invariant;
    \item
    \[
        s_1^\ast C\leq t_1^\ast C;
    \]
    \item
    \[
        C\leq\Box_1C;
    \]
    \item
    \[
        C=\Box_1C;
    \]
    \item for every \(n\geq0\) and \(0\leq i\leq n\),
    \[
        s_n^\ast C
        \leq
        (\nu_i^{(n)})^\ast C;
    \]
    \item every finite canonical residual execution beginning in \(C\)
    remains in \(C\) at every boundary.
\end{enumerate}
\end{theorem}

\begin{proof}
Residual invariance is exactly the first subobject inequality. By the
adjunction
\[
    s_1^\ast\dashv\forall_{s_1},
\]
this is equivalent to
\[
    C\leq\forall_{s_1}(t_1^\ast C)=\Box_1C.
\]
Identity executions imply \(\Box_1C\leq C\), so the inequality is
equivalent to equality.

The one-step condition propagates inductively along finite strings, giving
the all-boundary condition. Conversely, taking \(n=1\) and \(i=1\)
recovers the one-step condition. The final statement is the internal-language
form of the all-boundary subobject inequalities.
\end{proof}

\subsection{Generated closure as forward endpoint reachability}
\label{subsec:generated-closure-endpoint-reachability}

For
\[
    C\hookrightarrow\mathbf{Beh}_0,
\]
define
\[
    \operatorname{Reach}_{\mathrm{path}}(C)
    :=
    \operatorname{Im}
    \left(
        \coprod_{n\in\mathbb N}
        C\times_{\mathbf{Beh}_0,s_n}
        \mathscr T_{\mathbf{Beh}}(n)
        \xrightarrow{\ t_n\ }
        \mathbf{Beh}_0
    \right).
\]
Equivalently,
\[
    \operatorname{Reach}_{\mathrm{path}}(C)
    =
    \bigvee_{n\in\mathbb N}
    \operatorname{Post}_n(C).
\]

\begin{theorem}[Generated closure is endpoint reachability]
\label{thm:generated-closure-endpoint-reachability}
For every behavioural predicate \(C\),
\[
    \operatorname{Reach}_{\mathrm{path}}(C)
    =
    \operatorname{Post}_1(C)
    =
    \operatorname{Gen}_{\partial}(C).
\]
\end{theorem}

\begin{proof}
Total contraction gives
\[
    \operatorname{Post}_n(C)
    \leq
    \operatorname{Post}_1(C)
\]
for every \(n\). Since the one-step summand occurs in the defining
coproduct, the join of all finite endpoint images is
\(\operatorname{Post}_1(C)\).

Moreover,
\[
    C\times_{\mathbf{Beh}_0,s_1}
    \mathscr T_{\mathbf{Beh}}(1)
    \cong
    A_1\times_{t_A,A_0,p_C}C
    =
    D(C).
\]
Under this identification, the terminal map is exactly
\[
    \partial(1\times j_C):D(C)\to\mathbf{Beh}_0.
\]
Its regular image is the definition of
\(\operatorname{Gen}_{\partial}(C)\).
\end{proof}

Identity executions imply
\[
    C\leq\operatorname{Post}_1(C).
\]
Thus residual invariance is equivalently the fixed-point equation
\[
    \operatorname{Post}_1(C)=C.
\]

\subsection{The safety core as universal finite-path safety}
\label{subsec:safety-core-universal-finite-path-safety}

Define
\[
    \operatorname{Safe}_{\mathrm{path}}(C)
    :=
    \bigwedge_{n\in\mathbb N}
    \bigwedge_{0\leq i\leq n}
    \forall_{s_n}
    \left(
        (\nu_i^{(n)})^\ast C
    \right).
\]

The all-boundary formula is equivalent to the terminal-boundary formula
\[
    \operatorname{Safe}_{\mathrm{path}}(C)
    =
    \bigwedge_{n\in\mathbb N}
    \Box_nC.
\]
Indeed, the terminal conditions occur among the all-boundary conditions. For
the converse, put
\[
    S_i:=\forall_{s_i}(t_i^\ast C).
\]
The counit of
\[
    s_i^\ast\dashv\forall_{s_i}
\]
gives
\[
    s_i^\ast S_i\leq t_i^\ast C.
\]
Pulling back along the prefix restriction
\[
    r_{0,i}:
    \mathscr T_{\mathbf{Beh}}(n)
    \to
    \mathscr T_{\mathbf{Beh}}(i)
\]
gives
\[
    s_n^\ast S_i
    \leq
    (\nu_i^{(n)})^\ast C.
\]
Adjunction then gives
\[
    S_i
    \leq
    \forall_{s_n}
    \left(
        (\nu_i^{(n)})^\ast C
    \right).
\]
Taking the relevant meets proves the claim.

\begin{theorem}[Safety core as universal finite-path safety]
\label{thm:safety-core-universal-finite-path-safety}
For every behavioural predicate \(C\),
\[
    \operatorname{Safe}_{\mathrm{path}}(C)
    =
    \Box_1C
    =
    \operatorname{Safe}_{\partial}(C).
\]
Consequently, \(\operatorname{Safe}_{\partial}(C)\) is the largest subobject
from which every finite canonical residual execution remains in \(C\).
\end{theorem}

\begin{proof}
By
Proposition~\ref{prop:collapse-unconstrained-modalities},
\[
    \bigwedge_{n\in\mathbb N}\Box_nC
    =
    \Box_1C,
\]
where the length-zero factor is redundant because identity executions imply
\(\Box_1C\leq C\).

Under the identification
\[
    \mathscr T_{\mathbf{Beh}}(1)\cong D\mathbf{Beh},
\]
one has
\[
    s_1=\operatorname{pr}_\partial,
    \qquad
    t_1=\partial.
\]
Therefore
\[
    \Box_1C
    =
    \forall_{\operatorname{pr}_\partial}(\partial^\ast C),
\]
which is exactly the residual safety core
\(\operatorname{Safe}_{\partial}(C)\).
\end{proof}

\begin{corollary}[Execution meaning of the derivative adjoint triple]
\label{cor:execution-meaning-derivative-adjoint-triple}
The adjoint triple
\[
    \operatorname{Gen}_{\partial}
    \dashv
    I_{\partial}
    \dashv
    \operatorname{Safe}_{\partial}
\]
has the temporal interpretation:
\begin{enumerate}
    \item
    \[
        \operatorname{Gen}_{\partial}(C)
        =
        \operatorname{Post}_1(C)
    \]
    is forward residual reachability;
    \item \(I_\partial\) forgets residual invariance;
    \item
    \[
        \operatorname{Safe}_{\partial}(C)
        =
        \Box_1C
    \]
    is universal residual future safety.
\end{enumerate}
The residually invariant predicates are simultaneously the fixed points of
\(\operatorname{Post}_1\) and \(\Box_1\).
\end{corollary}

\section{One-object and Boolean execution semantics}
\label{sec:one-object-boolean-execution}

\subsection{One-object execution categories}
\label{subsec:one-object-execution-categories}

Let
\[
    \Phi:
    \mathbf{B}M\rightsquigarrow\mathbf{B}N
\]
be a one-object Mealy morphism. Recall that its carrier is an internal right
\(M\)-object \(P\), written
\[
    x\cdot m,
\]
together with an \(N\)-valued output cocycle
\[
    \omega:M\times P\to N.
\]

Using the cartesian symmetry, write execution arrows in state-first form
\[
    (x,m).
\]
Then
\[
    \operatorname{Exec}(\Phi)_0=P,
    \qquad
    \operatorname{Exec}(\Phi)_1=P\times M,
\]
with
\[
    s(x,m)=x,
    \qquad
    t(x,m)=x\cdot m,
\]
\[
    e(x)=(x,1_M),
\]
and
\[
    (x,m);(x\cdot m,n)
    =
    (x,mn).
\]

The input functor
\[
    \operatorname{Exec}(\Phi)\to(\mathbf{B}M)^{\mathrm{op}}
\]
is the discrete opfibration given by the action category, so
\(\operatorname{Exec}(\Phi)\) is the Grothendieck construction of the right
\(M\)-action.

An \(n\)-step execution is uniquely determined by
\[
    (x;m_1,\ldots,m_n).
\]
Its state path is
\[
\begin{tikzcd}[column sep=large]
    x
        \arrow[r, "m_1"]
    &
    x\cdot m_1
        \arrow[r, "m_2"]
    &
    \cdots
        \arrow[r, "m_n"]
    &
    x\cdot(m_1\cdots m_n).
\end{tikzcd}
\]
Total contraction sends it to
\[
    (x;m_1\cdots m_n).
\]

This displays concretely why path length is segmentation in the full
one-object category: an arbitrary monoid element already labels a single
categorical step.

\subsection{Boolean derivative execution}
\label{subsec:boolean-derivative-execution}

From this subsection onward, work in
\[
    \mathcal E=\mathbf{Set}.
\]
Let \(\Sigma\) be a set of letters and put
\[
    \mathbb A_\Sigma:=\mathbf{B}\Sigma^\ast,
    \qquad
    \mathbb B:=\mathbf 2_{\mathrm{ch}},
\]
where \(\mathbf 2_{\mathrm{ch}}\) is the chaotic category on
\[
    2=\{0,1\}.
\]

Let
\[
    L:\Sigma^\ast\to2
\]
be a Boolean language. Its left derivatives are
\[
    \partial_uL(w):=L(uw),
\]
and
\[
    \operatorname{Der}(L)
    :=
    \{
        \partial_uL
        \mid
        u\in\Sigma^\ast
    \}.
\]

Let
\[
    \operatorname{Min}(L):
    \mathbb A_\Sigma
    \rightsquigarrow
    \mathbf 2_{\mathrm{ch}}
\]
be the canonical derivative realization. Its transition is
\[
    H\cdot v:=\partial_vH,
\]
and its output anchor is
\[
    q(H):=H(\varepsilon).
\]
For a word \(v\), its output comparison is the unique arrow
\[
    \omega(v,H):
    H(v)=q(\partial_vH)
    \longrightarrow
    q(H)=H(\varepsilon)
\]
in \(\mathbf 2_{\mathrm{ch}}\).

Write
\[
    \operatorname{Exec}(L)
    :=
    \operatorname{Exec}\bigl(\operatorname{Min}(L)\bigr),
\]
and
\[
    \mathscr T_L
    :=
    \operatorname{Path}\bigl(\operatorname{Exec}(L)\bigr).
\]

The full derivative execution category has:
\begin{enumerate}
    \item objects \(H\in\operatorname{Der}(L)\);
    \item arrows
    \[
        (H,v):
        H\to\partial_vH,
        \qquad
        v\in\Sigma^\ast.
    \]
\end{enumerate}
Composition is
\[
\begin{tikzcd}[column sep=large]
    H
        \arrow[r, "v"]
    &
    \partial_vH
        \arrow[r, "w"]
    &
    \partial_{vw}H,
\end{tikzcd}
\]
contracting to \((H,vw)\).

\begin{proposition}[Finite word-labelled derivative executions]
\label{prop:finite-word-labelled-derivative-executions}
The \(n\)-sections of \(\mathscr T_L\) are tuples
\[
    (H;v_1,\ldots,v_n),
\]
where
\[
    H\in\operatorname{Der}(L),
    \qquad
    v_i\in\Sigma^\ast.
\]
The residual states are
\[
    H,
    \quad
    \partial_{v_1}H,
    \quad
    \partial_{v_2}\partial_{v_1}H,
    \quad
    \ldots,
    \quad
    \partial_{v_n}\cdots\partial_{v_1}H.
\]
If \(H=\partial_uL\), this becomes
\[
    \partial_uL,
    \quad
    \partial_{uv_1}L,
    \quad
    \partial_{uv_1v_2}L,
    \quad
    \ldots,
    \quad
    \partial_{uv_1\cdots v_n}L.
\]
\end{proposition}

\begin{proof}
Apply deterministic path lifting to the right
\(\Sigma^\ast\)-action on \(\operatorname{Der}(L)\).
\end{proof}

In
\(\mathbf 2_{\mathrm{ch}}^{\mathrm{op}}\), the output label of
\[
    (H,v):H\to\partial_vH
\]
is the unique arrow
\[
    H(\varepsilon)\longrightarrow H(v).
\]
Thus the output-labelled trajectory records the Boolean observation at every
state boundary.

\subsection{Letterwise execution paths}
\label{subsec:letterwise-execution-paths}

The full execution category treats every word as a one-step label. To make
path length count letters, retain only transitions labelled by elements of
\(\Sigma\).

\begin{definition}[Letterwise derivative graph]
\label{def:letterwise-derivative-graph}
The \textbf{letterwise derivative graph}
\[
    G_L^\Sigma
\]
has vertex set
\[
    \operatorname{Der}(L),
\]
edge set
\[
    \operatorname{Der}(L)\times\Sigma,
\]
and
\[
    s(H,a)=H,
    \qquad
    t(H,a)=\partial_aH.
\]
Its trajectory sheaf is
\[
    \mathscr T_L^\Sigma
    :=
    \operatorname{Path}(G_L^\Sigma).
\]
\end{definition}

An \(n\)-section is
\[
    (H;a_1,\ldots,a_n),
    \qquad
    a_i\in\Sigma.
\]
Starting at \(L\), its residual states are
\[
    L,
    \quad
    \partial_{a_1}L,
    \quad
    \ldots,
    \quad
    \partial_{a_1\cdots a_n}L.
\]
The observation at boundary \(i\) is
\[
    q\bigl(\partial_{a_1\cdots a_i}L\bigr)
    =
    L(a_1\cdots a_i).
\]

\begin{proposition}[Letterwise restriction]
\label{prop:letterwise-restriction}
For every \(n\geq0\), there is a canonical monomorphism
\[
    \mathscr T_L^\Sigma(n)
    \hookrightarrow
    \mathscr T_L(n)
\]
whose image consists of those word-labelled executions for which every
one-step word has length one. These monomorphisms commute with contiguous
restriction.
\end{proposition}

\begin{proof}
The inclusion
\[
    \Sigma\hookrightarrow\Sigma^\ast
\]
defines an inclusion of the letter-transition graph into the underlying
graph of the full word-action category. Apply the graph path functor.
\end{proof}

The letterwise trajectory sheaf is closed under restriction and gluing, but
not under contraction inherited from the full execution category:
contracting two letters produces a word of length two.

\subsection{Free categorical completion}
\label{subsec:free-categorical-completion}

Let
\[
    \operatorname{Free}(G_L^\Sigma)
\]
be the free category on the letterwise derivative graph.

\begin{theorem}[Word execution as free categorical completion]
\label{thm:word-execution-free-categorical-completion}
There is a canonical isomorphism
\[
    \operatorname{Free}(G_L^\Sigma)
    \cong
    \operatorname{Exec}(L).
\]
\end{theorem}

\begin{proof}
Define a functor that is the identity on residual languages and sends the
letter path
\[
    (H;a_1,\ldots,a_n)
\]
to
\[
    (H,a_1\cdots a_n).
\]
The empty path is sent to \((H,\varepsilon)\).

For a fixed \(H\), a letter sequence determines a unique path because every
successive residual state is forced. Every word has a unique decomposition
as a finite letter sequence. Hence the functor is bijective on hom-sets, and
concatenation of paths corresponds to multiplication of words.
\end{proof}

\begin{remark}[Letterwise paths are not already Segal closed]
\label{rem:letterwise-paths-not-segal-closed}
The full word-labelled nerve is not a Segal composition structure on the
fixed letterwise trajectory sheaf, since
\[
    \mathscr T_L^\Sigma(1)
    =
    \operatorname{Der}(L)\times\Sigma
\]
while
\[
    N_+\operatorname{Exec}(L)_1
    =
    \operatorname{Der}(L)\times\Sigma^\ast.
\]
A Segal composition structure leaves the degree-one object unchanged.
Rather, \(\operatorname{Exec}(L)\) is the free categorical completion of the
letterwise graph. It adjoins identity arrows labelled by the empty word and
composite one-step arrows labelled by arbitrary words.
\end{remark}

\subsection{Word deconcatenation as execution incidence coalgebra}
\label{subsec:word-deconcatenation-incidence}

A binary factorization of the execution arrow
\[
    (H,w):H\to\partial_wH
\]
is uniquely a factorization
\[
    w=uv.
\]
The intermediate state is \(\partial_uH\), and the corresponding two factors
are
\[
    (H,u):H\to\partial_uH,
\]
\[
    (\partial_uH,v):
    \partial_uH\to\partial_{uv}H.
\]

\begin{proposition}[Objective word deconcatenation]
\label{prop:objective-word-deconcatenation}
The incidence comultiplication of \(\operatorname{Exec}(L)\) sends
\((H,w)\) to the set of all cuts
\[
    w=uv
\]
together with the factor pair
\[
    \bigl((H,u),(\partial_uH,v)\bigr).
\]
\end{proposition}

If one applies a free-module linearization for which the relevant sums are
defined, the comultiplication is
\[
    \Delta(H,w)
    =
    \sum_{w=uv}
    (H,u)
    \otimes
    (\partial_uH,v).
\]
Coassociativity is the equality between the two ways of cutting a word into
three consecutive factors.

\subsection{Letterwise temporal modalities}
\label{subsec:letterwise-temporal-modalities}

For the letterwise trajectory sheaf, the endpoint modalities have genuine
length sensitivity:
\begin{enumerate}
    \item \(\Diamond_nC\) means that some word of length exactly \(n\)
    leads to a residual in \(C\);
    \item \(\Box_nC\) means that every word of length exactly \(n\) leads
    to a residual in \(C\);
    \item \(\operatorname{Post}_n(C)\) is exact \(n\)-letter forward
    reachability from \(C\).
\end{enumerate}

These operators do not collapse because the letterwise graph has no
one-step edge labelled by an arbitrary composite word. Passing to the free
categorical completion adds those composite arrows and thereby collapses the
unconstrained endpoint modalities to one categorical step.

\subsection{Classical finite-trajectory minimization}
\label{subsec:classical-finite-trajectory-minimization}

\begin{corollary}[Finite-trajectory minimization of deterministic automata]
\label{cor:finite-trajectory-minimization-deterministic-automata}
Let \(\mathcal A\) be a reachable deterministic automaton recognizing
\(L\). Its semantic quotient onto the derivative automaton induces:
\begin{enumerate}
    \item a surjection on executions of every finite path length for the full
    \(\Sigma^\ast\)-labelled semantics;
    \item a surjection on executions of every finite path length for the
    letterwise semantics;
    \item preservation of identity insertion and word contraction;
    \item a morphism of word-factorization incidence coalgebras;
    \item preservation of all letterwise temporal modalities under inverse
    image of state predicates.
\end{enumerate}
If \(\mathcal A\) is minimal, the trajectory maps are bijections and the
structured execution semantics are isomorphic.

For an arbitrary recognizing deterministic automaton, the same statements
hold after restriction to its reachable subautomaton.
\end{corollary}

\begin{proof}
For a reachable deterministic automaton recognizing \(L\), the semantic
state map is a surjective strict homomorphism onto
\(\operatorname{Min}(L)\). The structured trajectory minimization theorem
gives the full word-labelled statements.

For letterwise paths, deterministic lifting identifies a degree-\(n\)
execution with an initial state and an element of \(\Sigma^n\). The induced
map is the state quotient times the identity of \(\Sigma^n\), hence is
surjective. Naturality of inverse and direct image along this quotient gives
preservation of the modal formulas. If the automaton is minimal, its semantic
state map is an isomorphism.
\end{proof}

\section{Relation to interval-based system semantics}
\label{sec:relation-interval-based-semantics}

\subsection{Common span architecture}
\label{subsec:common-span-architecture}

Many sheaf-theoretic approaches to systems represent a machine by a span
\[
\begin{tikzcd}[column sep=large]
    \text{input trajectories}
    &
    \text{system trajectories}
        \arrow[l]
        \arrow[r]
    &
    \text{output trajectories}.
\end{tikzcd}
\]
The construction of this chapter has exactly this architecture:
\[
\begin{tikzcd}[column sep=large]
    \operatorname{Path}(\mathbb A^{\mathrm{op}})
    &
    \mathscr T_\Phi
        \arrow[l, "{\operatorname{Path}(I_\Phi)}"']
        \arrow[r, "{\operatorname{Path}(O_\Phi)}"]
    &
    \operatorname{Path}(\mathbb B^{\mathrm{op}}).
\end{tikzcd}
\]
Composition is pullback over the common intermediate trajectory sheaf:
\[
    \mathscr T_{\Psi\circ\Phi}
    \cong
    \mathscr T_\Phi
    \times_{\operatorname{Path}(\mathbb B^{\mathrm{op}})}
    \mathscr T_\Psi.
\]

The parallel is structural. In both cases, a system trajectory mediates
between compatible input and output data, and pullback synchronizes systems
along a shared interface.

\subsection{The four layers of temporal structure}
\label{subsec:four-layers-temporal-structure}

The constructions of the chapter distinguish four layers:
\begin{center}
\begin{tabular}{ll}
\textbf{Structure} & \textbf{Information retained} \\
\hline
Trajectory sheaf
& contiguous restriction and adjacent gluing \\
Labelled trajectory topos
& predicates, quantification, and synchronized interfaces \\
Segal structure
& identities, composition, and coherent contraction \\
Decomposition structure
& coherent factorization and incidence comultiplication.
\end{tabular}
\end{center}

A Mealy system is therefore not merely an arbitrary object of the labelled
trajectory topos. Its input projection is deterministic in the precise sense
that it is induced by an internal discrete opfibration, and its trajectory
object carries canonical Segal and decomposition structure.

\subsection{Derived rather than primitive temporal behaviour}
\label{subsec:derived-rather-than-primitive}

In an interval-based sheaf model, temporally extended behaviour is usually
specified directly. Sections over intervals and their restriction maps are
primitive, and gluing is part of the system definition.

Here the primitive data are instead:
\begin{enumerate}
    \item a persistent state carrier;
    \item typed input and output anchors;
    \item a deterministic transition operation;
    \item an output comparison satisfying the Mealy cocycle law.
\end{enumerate}
The trajectory sheaf is derived from the Grothendieck execution category. Its
gluing property follows from the iterated-pullback description of composable
one-step executions.

Every Mealy morphism therefore yields a discrete finite-time machine of
interval-sheaf form. The converse does not hold automatically: a general
object of the labelled trajectory topos need not arise from persistent state
and deterministic input lifting. The discrete-opfibration reconstruction
theorem gives the exact additional condition at the categorical execution
level.

The present interval site carries no metric duration or translation
structure. A section over \([n]\) records a path with \(n\) chosen
categorical segments. Sites with lengths, translations, continuous
durations, or hybrid time require a separate change-of-site comparison.

\subsection{The categorical and decomposition refinements}
\label{subsec:categorical-decomposition-refinements}

The finite-path sheaf records restriction and gluing. Since it comes from an
internal category, it also admits coherent identity insertion and
contraction:
\[
\begin{tikzcd}[column sep=large]
    N_+\operatorname{Exec}(\Phi)
        \arrow[r, "{j^\ast}"]
    &
    \mathscr T_\Phi.
\end{tikzcd}
\]

The same nerve also records coherent factorization. Its incidence
comultiplication sends a one-step execution to the object of all binary cuts.
Thus contraction and decomposition provide complementary views of
segmentation.

These refinements are essential for the residual semantic applications.
Contraction reduces finite-path endpoint reachability to one-step derivative
generation and universal finite-prefix safety to the one-step residual safety
core. Decomposition, by contrast, retains the intermediate states and all
ways of factoring a composite execution.

\subsection{Scope}
\label{subsec:temporal-scope}

The constructions of this chapter concern:
\begin{enumerate}
    \item discrete linear path shapes;
    \item finite executions;
    \item deterministic Mealy transitions;
    \item arbitrary Mealy simulation \(2\)-cells;
    \item sequential composition;
    \item strict Segal composition and objective execution decomposition;
    \item topos-internal temporal predicate logic.
\end{enumerate}
They do not by themselves provide:
\begin{enumerate}
    \item continuous-time or hybrid trajectories;
    \item infinitary execution objects;
    \item branching, cubical, pomset, or higher-dimensional execution shapes;
    \item a classification of all interval-based machines by Mealy
    morphisms;
    \item homotopy-coherent or higher-categorical execution semantics;
    \item a primitive metric notion of duration attached to a categorical
    arrow.
\end{enumerate}

A homotopical replacement by complete Segal objects and left fibrations is
left to a sequel. The next chapter instead studies a different locality
principle: variation of the input interface and reconstruction by descent
over interface covers.

\section{Chapter synthesis}
\label{sec:temporal-chapter-synthesis}

The chapter began from the execution category and output functor constructed
in the state-semantic development. The input label completed the two-sided
span
\[
\begin{tikzcd}[column sep=large]
    \mathbb A^{\mathrm{op}}
    &
    \operatorname{Exec}(\Phi)
        \arrow[l, "I_\Phi"']
        \arrow[r, "O_\Phi"]
    &
    \mathbb B^{\mathrm{op}}.
\end{tikzcd}
\]
The input leg is an internal discrete opfibration, and the execution category
is the internal Grothendieck construction of the state action. Conversely,
every such labelled discrete-opfibration span reconstructs a Mealy morphism.
Thus
\[
    \text{Mealy morphism}
    \quad\longleftrightarrow\quad
    \text{labelled deterministic execution category}.
\]

The construction retains the full simulation \(2\)-cells. A simulation
induces a functor between execution categories over the input interface and a
natural transformation comparing the output functors. This yields a
bicategorical execution representation rather than only a strict quotient
calculus.

Finite paths in the execution category are the finite runs of the Mealy
morphism. Deterministic lifting gives
\[
    \mathscr T_\Phi(n)
    \cong
    P\times_{A_0}
    \operatorname{Path}(\mathbb A^{\mathrm{op}})(n),
\]
so an execution is uniquely determined by its initial state and input path.
Pipeline composition is represented by pullback:
\[
    \mathscr T_{\Psi\circ\Phi}
    \cong
    \mathscr T_\Phi
    \times_{\operatorname{Path}(\mathbb B^{\mathrm{op}})}
    \mathscr T_\Psi.
\]

Abstracting contiguous restriction and gluing produced the equivalence
\[
    \mathbf{Graph}_{\mathrm{int}}(\mathcal E)
    \simeq
    \mathbf{TrajSh}_{\mathcal E}.
\]
Since internal graphs form the functor category
\(\mathcal E^{\mathbb G}\), trajectory sheaves form a Grothendieck topos.
Two-sided labelled trajectories form the slice topos
\[
    \mathbf{LabTraj}_{\mathbb A,\mathbb B}
    =
    \mathbf{TrajSh}_{\mathcal E}
    /
    \left(
        \operatorname{Path}(\mathbb A^{\mathrm{op}})
        \times
        \operatorname{Path}(\mathbb B^{\mathrm{op}})
    \right).
\]

Extending from contiguous interval maps to all simplicial maps produced the
equivalence
\[
    \mathsf{IntCat}(\mathcal E)
    \simeq
    \mathsf{Seg}^{\mathrm{str}}_+(\mathcal E).
\]
Thus graph paths carry restriction and gluing, while category paths
additionally carry coherent identities and contraction.

The same execution nerve is a decomposition object. It therefore determines
an objective incidence comonoid whose comultiplication records all binary
factorizations of an execution arrow. Mealy simulations induce CULF execution
functors and hence morphisms of these incidence comonoids. In the Boolean
word case, this comultiplication is word deconcatenation:
\[
    \Delta(H,w)
    =
    \sum_{w=uv}
    (H,u)\otimes(\partial_uH,v)
\]
whenever a suitable linearization is applied.

Regular semantic quotients remain regular epimorphic in every finite path
degree and preserve both Segal and factorization structure. Behavioural
minimization therefore induces simultaneous minimization of state-indexed
trajectories, their composition, and their execution incidence comonoids.

Finally, the internal logic of the trajectory topos supplies temporal
modalities. For canonical residual execution,
\[
    \operatorname{Gen}_{\partial}(C)
    =
    \operatorname{Post}_1(C)
\]
is forward residual reachability, while
\[
    \operatorname{Safe}_{\partial}(C)
    =
    \Box_1C
\]
is universal residual future safety. Hence
\[
    \text{residual invariance}
    \quad\Longleftrightarrow\quad
    \text{finite-prefix residual safety}
\]
and also
\[
    \text{residual invariance}
    \quad\Longleftrightarrow\quad
    \text{closure under forward residual reachability}.
\]

For full categorical execution, path length measures a chosen segmentation
rather than primitive duration. In the Boolean specialization, the
letterwise derivative graph supplies a primitive-step interpretation, the
full word-labelled execution category is its free categorical completion,
and the decomposition comonoid retains every possible cut of a word
execution.
